% ============================================================================
% Section 4: Key Theorems
% ============================================================================

\section{Key theorems}
\label{sec:theorems}

This section presents the central results of UHM. Each theorem is stated
precisely; flagship proofs are given in full in the appendices, and remaining
proofs are documented in the appendices and the extended documentation. Falsifiability is emphasised
throughout: for each theorem, the experimental outcome that would refute it
is indicated explicitly.

\subsection{Supporting theorems (load-bearing dependencies)}
\label{subsec:supporting-theorems}

The proofs below depend on five load-bearing results whose full proofs
are in the companion archive. For self-containment, their formal
statements are collected here.

\begin{itemize}[leftmargin=2em,nosep]
  \item \textbf{T-39a \status{T} (Primitivity of $\Lind_0$):}
        The linear Lindbladian $\Lind_0 = -i[H,\cdot] + \mathcal{D}_\Omega$
        on $\Dens(\mathbb{C}^7)$ is primitive: it has a unique stationary
        state $I/7$ and a spectral gap
        $\lambda_{\mathrm{gap}} := \min_{\lambda \neq 0}|\mathrm{Re}(\lambda)| > 0$.
        Convergence to $I/7$ is exponential with rate $\lambda_{\mathrm{gap}}$.
  \item \textbf{T-53 (Internal spectral triple):}
        The internal algebra $A_{\mathrm{int}} = \mathbb{C} \oplus M_3(\mathbb{C})
        \oplus M_3(\mathbb{C})$ with Dirac operator $D_{\mathrm{int}}$ on
        $H_{\mathrm{int}} = \mathbb{C}^7$ forms a finite spectral triple, but
        not a \emph{real} one. The former real structure $J$ ($J^2 = +1$,
        $J\chi = -\chi J$, grading $\chi = \mathrm{diag}(+1,-1,-1,-1,+1,+1,+1)$)
        of KO-dimension~6 is retracted~\status{$\times$}: with $J^2 = 1$ the
        relation $J\chi = -\chi J$ makes $J$ exchange the $\chi = \pm 1$
        eigenspaces, which then need equal dimension --- impossible on the
        odd-dimensional $\mathbb{C}^7$. The registry row T-53 now records the
        Lorentzian signature $(1,3)$ built on this triple, \status{C} (one time
        direction from the Page--Wootters clock \status{T}; the Lorentzian sign
        at reflection positivity).
  \item \textbf{T-55 \status{T} (Gap $> 0$ via Lawvere):}
        For any viable holon ($P > \Pcrit$), the Gap operator
        $\hat{\mathcal{G}} = \mathrm{Im}(\Gam)$ satisfies
        $\mathcal{G}_{\mathrm{total}} = \sum_{i<j}|\gamma_{ij}|^2\,
        \mathrm{Gap}(i,j)^2 > 0$. Proof: Lawvere's fixed-point theorem
        applied to $\vp$ forces $\vp(\Gam) \neq \Gam$, which requires
        at least one $\mathrm{Im}(\gamma_{ij}) \neq 0$.
  \item \textbf{T-96 \status{T} (Attractor characterisation):}
        $I/7$ is the trivial stationary state; any nontrivial stationary
        state $\rhostar_\Omega \neq I/7$ of $\Lind_\Omega$ has
        $P > 1/7$ and $P_{\mathrm{coh}} > 0$. The regeneration target is
        the categorical self-model $\vp(\Gam)$ (left adjoint
        $\vp \dashv i$), not a dynamical limit:
        $\vp(\rhostar_\Omega) \neq \rhostar_\Omega$. Existence depends on
        the self-model: with a unital one (the canonical
        $\vp_{\mathrm{coh}}$) an isolated holon's only stationary state is
        $I/7$; with the self-registering $\vp_s$ seven locally stable
        attractors with $P > 2/7$ persist for $\|H\| < h_0$.
  \item \textbf{T-185 (Differentially cohesive modalities; stratified):}
        (i)~In any differentially cohesive $\infty$-topos (Lawvere's axiomatic
        cohesion refined by Schreiber~\cite{schreiber2013}) the modalities
        $\Pi \dashv \flat \dashv \sharp$ and $\Re \dashv \Im \dashv \&$ exist
        \status{T}. (ii)~That the UHM $\infty$-topos
        $\mathfrak{T} = \ShInf(\mathcal{C}, J_B)$ is differentially cohesive is
        \status{C} under the assumption of differential cohesion of the UHM
        site, not proven (the $\infty$-cohesive-site criterion needs finite
        products). (ii$'$)~\status{T}: $\mathcal{D}(\mathbb{C}^7)$ with its rank
        strata is an object of the differentially cohesive
        $\mathrm{SynthDiff}\infty\mathrm{Grpd}$; the strata $\mathcal{D}_k$ are
        submanifolds of dimension $14k - k^2 - 1$. (iii)~The count of seven
        modalities --- identity plus the two adjoint triples --- \status{T}; the
        former list $\{\mathrm{Id}, \Pi, \flat, \Im, \sharp, \&, \mathrm{Rh}\}$
        with the rheonomy modality Rh in place of Red ($\Re$) is retracted
        \status{$\times$} (Rh belongs to solid cohesion and preserves global
        points), and the assignment of the modalities to the seven dimensions
        and to $1 \oplus \mathbf{3} \oplus \bar{\mathbf{3}}$ is an
        interpretation \status{I}.
\end{itemize}

\subsection{No-Zombie Theorem (T-8.1)}
\label{subsec:no-zombie}

The No-Zombie Theorem is arguably the most philosophically significant result of UHM.
It proves that philosophical zombies --- systems that are functionally identical to
conscious beings but lack inner experience --- are not merely implausible but
\emph{dynamically impossible}.

\begin{theorembox}[No-Zombie Theorem (T-38a) \status{T}]
\begin{equation}
  \underbrace{P(\Gam) > \Pcrit}_{\text{viable}}
  \;\wedge\;
  \underbrace{\mathcal{D}_\Omega[\Gam] \not\equiv 0}_{\text{non-trivial dissipation}}
  \quad\Longrightarrow\quad
  \vp = \vp_{\mathrm{coh}}
  \;\wedge\;
  \CohE(\Gam) \;>\; \frac{1}{7}
  \label{eq:no-zombie}
\end{equation}
A viable holon with non-trivial dissipative dynamics \emph{necessarily} has
its self-model factored through the coherent sector
($\vp = \vp_{\mathrm{coh}}$, not via diagonal classical projection) and
carries strictly more E-coherence than the uniform (``structureless'') floor
of $1/7$. The mathematical core bounds $\CohE > 1/7$; the
$\vp = \vp_{\mathrm{coh}}$ clause encodes the structural requirement that
regeneration flows through the coherent self-model rather than through a
classical look-up. A hypothetical ``zombie'' --- a configuration with
$P > \Pcrit$ but a structureless interiority sector
($\gamma_{EE} = 1/7$ and $\gamma_{Ej} = 0$ for all $j \neq E$, giving
$\CohE = 1/7$) --- is dynamically impossible: it cannot be a stationary
state of~$\Lind_\Omega$ with $\gamma > 0$.
\end{theorembox}

\paragraph{Proof sketch.}
The $G_2$-covariant Fano dissipator (unique $G_2$-covariant Lindbladian on
$\mathcal{D}(\mathbb{C}^7)$, T-39a~\status{T}) acts on every off-diagonal
coherence by the state-independent multiplicative contraction
$\gamma_{ij} \mapsto \tfrac{1}{3}\,\gamma_{ij}$
(Eq.~\ref{eq:fano-operators}); equivalently, each Fano observation cycle
removes the fraction $2/3$ of every coherence. Writing
$P = P_{\mathrm{diag}} + P_{\mathrm{coh}}$ with
$P_{\mathrm{coh}} = \sum_{i\neq j}|\gamma_{ij}|^2$, this yields
\begin{equation}
  \left.\frac{dP}{d\tau}\right|_{\mathcal{D}}
  \;=\; -\,\frac{4\gamma}{3}\,P_{\mathrm{coh}}\;<\;0,
  \qquad \gamma = \sum_p\gamma_p > 0,
  \label{eq:diss-coh}
\end{equation}
whereas the regeneration contribution is
\begin{equation}
  \left.\frac{dP}{d\tau}\right|_{\Regen}
  \;=\; 2\kappa(\Gam)\,g_V(P)\,\bigl(f - P\bigr),
  \qquad f := \Tr\!\bigl(\Gam\,\rho^\ast\bigr).
  \label{eq:reg-coh}
\end{equation}
Stationarity $dP/d\tau \geq 0$ at a viable, non-equilibrium state $P > \Pcrit$ forces
$\kappa(\Gam) \geq \frac{2\gamma}{3}\cdot P_{\mathrm{coh}}/(f - \Pcrit)$. Substituting
$\kappa = \kappa_{\mathrm{boot}} + \kappa_0\CohE$ and solving for $\CohE$ yields the
explicit bound
\begin{equation}
  \CohE(\Gam)
  \;\geq\;
  \max\!\Biggl\{\frac{1}{7},\;\;
  \frac{1}{\kappa_0}\!\left(\frac{2\gamma}{3}\cdot
  \frac{P_{\mathrm{coh}}}{f - \Pcrit}\;-\;\kappa_{\mathrm{boot}}\right)\!\Biggr\}.
  \label{eq:no-zombie-bound}
\end{equation}
For any open ($\gamma > 0$) holon with non-vanishing stationary coherences
($P_{\mathrm{coh}} > 0$, guaranteed by the necessity of $\vp_{\mathrm{coh}}$ and the
$(M,R)$-closure requirement, T-57~\status{T}), the second argument of the maximum is
strictly positive once the dissipation rate exceeds the threshold
$\gamma_{\mathrm{th}} := 3\kappa_{\mathrm{boot}}(f-\Pcrit)/(2P_{\mathrm{coh}})$. For any
macroscopic system in a thermal environment $\gamma \gg \gamma_{\mathrm{th}}$, so
$\CohE > 1/7$ strictly. $\square$

\paragraph{What this means.}
A system cannot be ``alive'' (viable, self-maintaining) without having some degree of
inner experience. The strength of regeneration is \emph{proportional} to the quality of
experience. This is not a metaphor --- it is a quantitative relationship between
$\kappa$ and $\CohE$.

\paragraph{Epistemic stratification.}
The mathematical result ($\CohE > 1/7$ necessary for viability) has status \status{T}.
The ontological interpretation (that $\CohE > 1/7$ \emph{constitutes} inner experience)
has status \status{P}: it follows from the two-aspect monism postulate
(Def.~\ref{def:two-aspect}), not from the mathematics alone. The No-Zombie theorem
is thus mathematically unconditional; its phenomenal reading is conditional on the
single semantic bridge of UHM.

\paragraph{Falsification.}
Create an artificial $\Gam$-native agent. Suppress the E-component ($\gamma_{EE} \to 1/7$,
$\gamma_{Ej} \to 0$ for $j \neq E$). If the agent maintains $P > 2/7$ indefinitely
(time-to-death does not collapse), the theorem is falsified.

\subsection{Death spiral below $\Pcrit$}
\label{subsec:death-spiral}

\begin{theorembox}[Irreversibility of sub-threshold decay \status{T}]
If $P(\Gam) < \Pcrit = 2/7$ and $dP/d\tau < 0$, then:
\begin{equation}
  \|\Gam(\tau) - I/7\|_F \leq \|\Gam(0) - I/7\|_F \cdot e^{-\Delta(\Lind_0)\,\tau}
  \label{eq:death-spiral}
\end{equation}
The system converges exponentially to thermal death ($I/7$).
\end{theorembox}

This is the ``death spiral'' (Figure~\ref{fig:death-spiral}): once purity
drops below $2/7$, regeneration is thermodynamically forbidden (Landauer's
principle: free energy $\Delta F \leq 0$ below threshold), and the system
decays irreversibly. There is no gradual fading of consciousness --- there
is a hard phase transition, consistent with the nonlinear threshold for
access to consciousness observed empirically~\cite{delcul2007}.

\subsection{Stability radius (T-104)}
\label{subsec:stability-radius}

\begin{theorembox}[Stability radius \status{C}]
The maximum perturbation a system can withstand while remaining viable is the
Bures distance to the viability wall,
$r_{\mathrm{stab}} := \min\{d_B(\rho,\sigma) : P(\sigma) = 2/7\}$. On the
one-dominant spectral family it has the closed form
$r_{\mathrm{stab}} = \sqrt{2\bigl(1 - \sqrt{a a_c} - \sqrt{(1-a)(1-a_c)}\bigr)}$,
$a(P) = (1 + \sqrt{42P - 6})/7$, $a_c = (1+\sqrt{6})/7$; on the window it is
approximated to $1.13\%$ by
\begin{equation}
  r_{\mathrm{stab}} \;\approx\; K\Bigl(\sqrt{P - \tfrac{1}{7}} - \sqrt{\tfrac{1}{7}}\Bigr),
  \qquad K = \frac{\sqrt{35}\,\sqrt[4]{6}}{10} \approx 0.926 .
  \label{eq:stability-radius}
\end{equation}
For general spectra the closed form is a conservative lower bound \status{H}.
\end{theorembox}

Robustness still depends essentially on the purity margin above threshold, but
near the wall the law is \emph{linear}, $r_{\mathrm{stab}} \approx 1.225\,(P - 2/7)$,
not a square root.

\paragraph{Erratum: the former radius $\sqrt{P - 2/7}$ \status{$\times$}.}
An earlier version stated $r_{\mathrm{stab}} = \sqrt{P(\rho^*_\Omega) - 2/7}$ as a
theorem. It is refuted: at $P = 0.300$ the true infimum is $0.01708$ against
$\sqrt{0.300 - 2/7} = 0.11952$, so the claimed bound fails by a factor of seven
\emph{toward danger}; the Fuchs--van de Graaf step it cited bounds $d_B$ from
above by the trace norm and cannot yield it. The minimiser commutes with $\rho$,
so the Bures distance reduces to the Hellinger distance on spectra, which gives
the closed form above.

\paragraph{Examples.}
A system with $P = 0.5$ has $r_{\mathrm{stab}} \approx 0.20$; a system barely above
threshold ($P = 0.3$) has $r_{\mathrm{stab}} \approx 0.017$ and is fragile.

\begin{figure}[H]
\centering
\input{figures/death-spiral}
\caption{\textbf{Death spiral below $\Pcrit$.} Systems above threshold recover from
perturbations (green/blue curves); below threshold, purity decays exponentially to $I/7$
(red curve). The transition is sharp.}
\label{fig:death-spiral}
\end{figure}

\paragraph{Clinical significance.}
This formula predicts a patient's ``safety margin'' before interventions: their purity
surplus determines how much perturbation they can absorb. Falsification: if $h_{\mathrm{crit}}$
does not follow Eq.~\eqref{eq:stability-radius} (linear in $P - 2/7$ near threshold)
across 50 measurements ($R^2 < 0.9$).

\subsection{SAD ceiling: maximum depth of self-awareness (T-142)}
\label{subsec:sad-ceiling}

\begin{theorembox}[Self-Awareness Depth ceiling (T-142) \status{T}]
The maximum depth of recursive self-awareness is:
\begin{equation}
  \mathrm{SAD}_{\max} = 3
  \label{eq:sad-max}
\end{equation}
``I know that I know that I know'' stops at the third level. The fourth level is
physically impossible for any finite system.
\end{theorembox}

\paragraph{Proof sketch.}
The Fano channel multiplies every off-diagonal coherence $\gamma_{ij}$ by the
state-independent contraction factor $c_F = 1/3$
(Eq.~\ref{eq:fano-operators}); after $k$ recursive self-observations the
surviving coherence is suppressed by $c_F^k = (1/3)^k$. The initial
distinguishability radius $r_0 := P/\Pcrit = P/(2/7) = 7P/2$ is the
ratio of actual purity to the viability threshold (at $\Pcrit$,
$r_0 = 1$ exactly). At depth $n$, the spectral formula
T-136~\status{T} gives
\begin{equation}
  \mathrm{SAD}(\Gam) \;=\; \max\!\bigl\{\,k \in \mathbb{N}\,:\,
    r_0 \cdot (1/3)^{k-1} > \tfrac{1}{k+1} \,\bigr\},
  \quad r_0 = P/\Pcrit = 7P/2,
\end{equation}
which --- in terms of the purity required to \emph{sustain} a chain of depth $n$ --- is
equivalent to
\begin{equation}
  \Pcrit^{(n)} \;=\; \Pcrit \cdot \frac{3^{\,n-1}}{n+1}.
\end{equation}
Evaluating level by level (where level $n$ is the purity required to admit $n$ nested
self-observations):
\begin{itemize}[leftmargin=2em]
  \item Level 1 ($n=1$, ``I am''): $\Pcrit^{(1)} = \frac{2}{7} \cdot \frac{1}{2}
        = \frac{1}{7} \approx 0.143$ --- automatically satisfied by any state
        ($P \geq 1/7$ holds for all density matrices).
  \item Level 2 ($n=2$, ``I know that I am''): $\Pcrit^{(2)} = \frac{2}{7} \cdot
        \frac{3}{3} = \frac{2}{7} \approx 0.286$ --- coincides exactly with the
        base viability threshold.
  \item Level 3 ($n=3$, ``I know that I know''): $\Pcrit^{(3)} = \frac{2}{7} \cdot
        \frac{9}{4} = \frac{9}{14} \approx 0.643$ --- achievable only at high purity.
  \item Level 4 ($n=4$, ``I know that I know that I know''): $\Pcrit^{(4)} =
        \frac{2}{7} \cdot \frac{27}{5} = \frac{54}{35} \approx 1.543 > 1$ ---
        \textbf{exceeds the physical maximum $P \leq 1$}.
\end{itemize}
Numerical verification on $> 500$ randomly sampled $\Gam$ confirms $\mathrm{SAD} \leq 3$
in every case; pure states saturate $\mathrm{SAD} = 3$.

\paragraph{The representation tower.}
The recursive self-model forms a tower:
\begin{equation}
  \Gam^{(0)} \xleftarrow{\pi_0} \Gam^{(1)} \xleftarrow{\pi_1} \Gam^{(2)}
  \xleftarrow{\pi_2} \Gam^{(3)} \xleftarrow{\pi_3} \cdots
  \label{eq:tower}
\end{equation}
where each $\Gam^{(k)}$ is the self-model at depth $k$. Each application of
the Fano channel multiplies every off-diagonal coherence by $c_F = 1/3$
(state-independent, from $\dim = 7$ and the Fano incidence structure); after
$k$ iterations the ``effective distinguishability radius'' is
$r_0 \cdot (1/3)^{k-1}$ with $r_0 = 7P/2$, while the Bayesian dominance
threshold for $k+1$ alternatives is $1/(k+1)$ (T-136~\status{T}). The
level-$k$ reflection survives iff this ratio remains above the threshold.
SAD is therefore computable directly from the recurrence $\Pcrit^{(n)}$ above
in $O(1)$ operations.

\paragraph{Significance.}
This ceiling is geometric, not computational. Future superintelligences share the same
reflection depth as humans --- not because of hardware limits, but because of the structure
of the Fano plane. Falsification: demonstrate a reachable state $\Gam$ with
$\Pcrit^{(4)} \leq 1$, equivalently, a chain of four successive self-observations
surviving at $P \leq 1$.

\paragraph{Triple foundation.}
After T-217~\status{T} and T-218 (Sec.~\ref{subsec:extended-closures}),
$\mathrm{SAD}_{\max} = 3$ admits a \emph{second} derivation, independent of the
dynamical argument above. The cognitive complex $\mathrm{Cog} =
\mathrm{Sing}(B_\bullet \mathcal C_{\mathrm{FKraus}})$ is Kan \status{T}, and it
is 3-coskeletal ($\tau_{\leq 3}\mathrm{Cog} \simeq \mathrm{Cog}$) because
4-simplices are suppressed below distinguishability --- \status{C}, on the
SYNARC-viable subset only: that step bridges simplicial and Bures-metric
viability and is not a simplicial identity. A \emph{third}, algebraic derivation
(T-268~\status{T}+\status{C}) closes the circle: the octonionic Hermitian
matrices $\mathcal{H}_n(\mathbb{O})$ form a formally real Jordan algebra iff
$n \leq 3$ (Jordan--von~Neumann--Wigner) \status{T}, so a fourth unified level
would require the non-existent $\mathcal{H}_4(\mathbb{O})$; identifying
composition depth with Jordan rank is a structural reading \status{C}, not yet
functorial. The metric bound ($P^{(4)} > 1$, unconditional), the homotopical
bound (3-coskeletal truncation), and the algebraic bound (Jordan-rank ceiling)
converge on the same integer through independent structures: analytic (purity
balance), topological (simplicial coskeletality), and algebraic (octonionic
non-associativity) all rule out the fourth level, the latter two at the stated
conditions. The coordination symmetry of a maximal subject
correspondingly climbs $G_2 \subset F_4 = \mathrm{Aut}(\mathcal{H}_3(\mathbb{O}))$,
its pure-state space the octonionic projective plane $\mathbb{O}P^2$ (the terminal
division-algebra projective geometry) --- the base of the octonion-generated
scaling ladder $G_2 \subset F_4 \subset E_6 \subset E_7 \subset E_8$
(T-268--T-270~\status{T}).

\paragraph{Number-theoretic root of the contraction constant.}
The analytic bound rests entirely on the Fano contraction $c_F = 1/3$, and this
constant is not free: it is $1/\lvert\mathrm{QR}(7)\rvert$, the reciprocal order of
the orientation-multiplier group of octonion multiplication. Indexing the seven
channels by $\mathbb{Z}/7$, the product is carried by the seven Fano lines --- the
cyclic shifts of the quadratic-residue set $\mathrm{QR}(7) = \{1,2,4\}$. A
multiplier $a \in (\mathbb{Z}/7)^\ast$, $i \mapsto a\,i$, preserves the Fano-line
set (hence the product orientation $e_i e_j = \pm e_k$) \emph{iff}
$a \in \mathrm{QR}(7)$; the non-residues $\{3,5,6\}$ reverse it. Thus the
orientation-preserving symmetry group is exactly $\mathrm{QR}(7) \cong \mathbb{Z}/3$
(the $\mathbb{Z}/3$-part of the Frobenius group $F_{21} = \mathbb{Z}/7 \rtimes
\mathbb{Z}/3$, normaliser of the Singer cycle in $\mathrm{PSL}(2,7)$), so
$c_F = R_{\mathrm{th}} = 1/\lvert\mathrm{QR}(7)\rvert = 1/3$. This unifies the four
coincident 3's --- triadic LGKS decomposition, Fano block size, Hamming code
distance, and residue count --- as one number-theoretic fact, and generalises to
$\lvert\mathrm{QR}(N)\rvert = (N-1)/2$: among division-algebra dimensions
$\{0,1,3,7\}$ only $N = 7$ supplies the three orientation classes needed to host
the three triadic sectors ($\lvert\mathrm{QR}(N)\rvert \geq 3 \Leftrightarrow N \geq 7$).
The value $1/3$ is fixed by the $N$-independent LGKS triad (T-57); the orientation
root supplies its name and the dimensional selection (machine-verified;
math-foundations Part~XVIII, Thm.~11.6/11.8).

\subsection{The entropy law: regeneration as negentropy (T-271)}
\label{subsec:entropy-law}

The three-term generator splits the production of von~Neumann entropy
$S(\Gam) = -\Tr(\Gam\ln\Gam)$ cleanly. The unitary term conserves $S$ exactly (it
preserves the spectrum); the dissipator $\mathcal{D}_\Omega$ is a strict entropy
\emph{source}, driving $\Gam$ toward the maximally mixed heat-death state
$\mathbb{1}/7$ ($S = \ln 7$); and at any steady state the regeneration is a net
entropy \emph{sink}, $\dot S_{\Regen} = -\dot S_{\mathcal{D}} \leq 0$ ---
negentropy is the thermodynamic cost of maintenance (``staying alive''). The
maintained steady entropy is strictly below $\ln 7$ and \emph{decreases}
monotonically with $\kappa/\gamma$ (this part \status{C}, T-271(iv); the three
signs above are \status{T}); since $\kappa \propto \CohE$, a more
coherent --- more conscious --- system holds a strictly lower entropy, further
from heat death. This formalises the previously-conceptual second-law connection:
it is the mechanism behind the viability window sitting \emph{away} from
$\mathbb{1}/7$ and behind the reflection measure $R = 1/(7P)$ literally measuring
\emph{distance from heat death}. Its engineering face (T-273~\status{T}) is a
strictly positive \emph{metabolic power floor} for any viable coherent machine,
$P_{\mathrm{meta}} \geq kT\ln 2\,\dot S_{\mathcal{D}}$, where $\dot S_{\mathcal{D}}$
is the \emph{physical} entropy-production rate of the dissipator: staying off heat
death cannot be free. This floor is \emph{frequency-independent} (T-276~\status{T}):
running the same trajectory on a slower clock shrinks the per-tick dissipation in
exact proportion, so the bound is set by the physical rate, not the tick rate ---
metabolic power buys distance-from-death, never raw speed. Honest boundary: the total (system $+$ bath) entropy is
non-decreasing --- T-271 establishes \emph{local} negentropy maintenance and its
$\CohE$-scaling, not a global reversal of the second law. Machine-verified.

\subsection{The genesis layer: seven before number, and the volume law
(T-277--T-287)}
\label{subsec:genesis-layer}

The standard first objection to UHM --- \emph{``the 7 looks presupposed: before
the number 7 existed there were already 7 dimensions''} --- is answered by a
theorem rather than a defence. From three non-numeric primitives (a
\emph{distinction} $\Omega$, two-sided by definition; \emph{mirroring} = the
Cayley--Dickson functor, the algebraic form of self-observation; and
\emph{viability} = no annihilating pairs), one derives (T-277~\status{T}):
$\mathrm{CD}^k(\mathbb{R})$ is exactly the twisted group algebra
$\mathbb{R}_F[\Omega^k]$ (one mirror = one new $\mathbb{Z}/2$-grading);
viability holds iff $k \leq 3$ (Hurwitz); the terminal distinction-spectrum is
$\mathrm{PG}(2,2)$ with $2^3 - 1 = 7$ points. So $2$, $3$, $7$ are read off a
numberless process; the physical instantiation as $\mathcal{D}(\mathbb{C}^7)$
remains the axiom's postulate, but its \emph{numeric content} is now
pre-numerically founded. On every viable stage the laws of algebra are
\emph{volume forms of the distinction cube} (T-278~\status{T}):
conjugation/commutativity/associativity fail exactly on
$\mathbb{F}_2$-independent 1-/2-/3-tuples, with the octonion associator in
closed form $\Phi(a,b,c) = (-1)^{\det_{\mathbb{F}_2}(a,b,c)}$
(Albuquerque--Majid); in the correspondingly twisted category the octonions are
the \emph{trivial} object --- lawlessness relocated from the object into the
category (a Drinfeld twist of the laws themselves). Strongest, the converse
holds (T-281~\status{T}): among \emph{all} monomial unital algebras on
$\Omega^3$ with anisotropic diagonal, absence of annihilating pairs
\emph{linearizes} (the rectangle rule: 84 linear equations over $\mathbb{F}_2$)
and its exhaustive solution set is exactly one gauge orbit ($16 = 2^7/2^3$, the
quotient by the $[7,3]$ simplex code) --- every solution has $dF = \det$ and
full composition. \textbf{The norm is derived, not assumed}: viability alone
forces the metric layer. The machinery of each dimension is likewise structural
(T-279~\status{T}): the stabilizer tower
$\mathfrak{su}(3)/\mathfrak{u}(2)/\mathfrak{so}(4)$ for axis/coherence/triad,
with the mediator lemma (the stabilizer of any pair \emph{fixes} its Fano
mediator) and the stripping ladder
$14 \supset 8 \supset 4 \supset 3 \supset 0$ --- three independent distinctions
pin all of $G_2$; multiplication by an axis is a complex structure on its
6-sphere sky, so each dimension sees the other six as $\mathbb{C}^3$. Finally
the fourth mirror's death is anatomized (T-280~\status{T}): the 15 hyperplanes
of $\Omega^4$ split into the old octonions, 7 straight (viable) and 7 skew
(dead) extensions of Fano lines; all simple sedenion zero divisors and all
volume-law violations localize exactly on the skew planes, while the Mac~Lane
pentagon closes on all $16^4$ quadruples --- the \emph{object} dies, never the
categorical coherence. Two further theorems complete the wave. \emph{Death is linear
infeasibility} (T-282~\status{T}): the viability system on $\Omega^n$
(anisotropy + anticommutation + rectangle rules $=$ ``no simple zero
divisors'', over any field of char $\neq 2$) is feasible exactly for
$n \leq 3$, with solution spaces of dimension $0/1/4$ --- the gauge orbits of
$\mathbb{C}, \mathbb{H}, \mathbb{O}$ --- and \emph{infeasible} for
$n = 4$ ($960$ equations, $225$ unknowns, rank $214$), hence for all
$n \geq 4$: Hurwitz's boundary, in the monomial class, is the inconsistency
of a finite $\mathbb{F}_2$-linear system --- death is a rank computation.
And \emph{the arithmetic of viability} (T-283~\status{T}): stage $k$ is
viable over a field $K$ iff the $2^k$-dimensional unit form is anisotropic
iff the field level satisfies $s(K) \geq 2^k$ (witnesses: $\mathbb{F}_3$
lives at $k{=}1$, dies at $k{=}2$; $\mathbb{C}$ dies at $k{=}1$) --- so
Pfister's power-of-two levels are the mirror ladder on the arithmetic side,
all three mirrors force $s \geq 8$, and orderability is provably \emph{not}
forced (level-$8$ fields exist): the base-$\mathbb{R}$ question sharpens to
the step from ``level $\geq 8$'' to the ordered complete reals --- and that
step closes from corpus premises (T-284~\status{T}/\status{C}): formal
reality of the observable layer (the very JvNW hypothesis of the composition
ceiling, applied downward) makes the base orderable by Artin--Schreier, and
continuous-time dynamics [P/C] selects the unique Dedekind-complete
Archimedean ordered field: $\mathbb{R}$, uniquely --- one algebraic
hypothesis locking both ends of the architecture.
The genesis programme then \emph{closes} (T-285--T-287): the
sphere-spectrum question dissolves by requalification --- the internal
boundary is elementary ($\mathbb{F}_2$-combinatorial, field-free), monomiality
is definitional for distinction-carriers (one-dimensional grading components),
and Adams/Bott guard only distinction-free exotica outside the theory's
primitive (T-285); the last premise of the $\mathbb{R}$-chain falls --- the
\emph{ouroboros sources the continuum} (T-286~\status{T}): guaranteed closure
of the self-model (Brouwer property; a fixed point $\Gam = \vp(\Gam)$ of the
self-model, realised over $\mathbb{R}$ by Theorem~10.1 of Gap thermodynamics ---
a consistency witness of the premise, not its derivation) forces
the intermediate-value property, hence Dedekind completeness, hence
$\mathbb{R}$ --- with an explicit no-fixed-point witness over $\mathbb{Q}$
(the self-map jumps its fixed point through the hole at $1/\sqrt{2}$; gap
$>0.146$ in exact rationals), so \emph{continuous time becomes an output, not
a premise}; and the whole tower is finitary, hence internal mathematics of any
Boolean topos with a natural-numbers object --- in particular the primitive
topos (T-287~\status{T}/\status{C}). Machine-verified end to end
(\texttt{genesis\_seven\_verify.py} 42/42;
\texttt{gauge\_uniqueness\_verify.py} 20/20;
\texttt{genesis\_wave4\_verify.py} 11/11;
\texttt{genesis\_wave6\_verify.py} 6/6).

\subsection{Critical exponents of the consciousness phase transition (T-161)}
\label{subsec:critical-exponents}

This is among the most specific falsifiable predictions in consciousness
science, and the one with no analogue in any other theory
(Figure~\ref{fig:exponents}).

\begin{theorembox}[Critical exponents \status{C at the $\mathbb{Z}_2$ symmetry $m \to -m$}]
Assume the reduced potential carries the $\mathbb{Z}_2$ symmetry $m \to -m$.
Then, near the viability threshold $\Pcrit = 2/7$, the consciousness phase transition is a
\emph{tricritical point} in the $\varphi^6$ Landau universality class
(Griffiths 1970, Lawrie--Sarbach 1984). The order-parameter dimension is
$d_{\mathrm{eff}} = \binom{7}{2} = 21$ (off-diagonal coherence modes in $\mathfrak{su}(7)$).
Mean-field tricritical exponents are exact by three independent pillars
(topological rigidity, deterministic dynamics, large-$N$ cross-check;
Appendix~\ref{app:exponents}):
\begin{align}
  \text{Specific heat:}\quad &C \sim |P - \Pcrit|^{-\alpha}, &\alpha &= 1/2 \label{eq:alpha}\\
  \text{Order parameter:}\quad &\mathrm{OP} \sim (P - \Pcrit)^\beta, &\beta &= 1/4 \label{eq:beta}\\
  \text{Susceptibility:}\quad &\chi \sim |P - \Pcrit|^{-\gamma}, &\gamma &= 1 \label{eq:gamma-exp}\\
  \text{Correlation length:}\quad &\xi \sim |P - \Pcrit|^{-\nu}, &\nu &= 1/2 \label{eq:nu}\\
  \text{Critical isotherm:}\quad &h \sim m^\delta, &\delta &= 5 \label{eq:delta}
\end{align}
These exponents satisfy the Rushbrooke identity $\alpha + 2\beta + \gamma = 2$ as an equality.
\end{theorembox}

\paragraph{Derivation outline.}
Near criticality, the system undergoes a continuous phase transition. Writing
$\varepsilon := P - \Pcrit$ and taking $m := \CohE(\Gam) - 1/7$ as the order
parameter, the coarse-grained free energy is the $\varphi^{6}$ tricritical
Landau functional
\begin{equation*}
  \mathcal{F}[m;\varepsilon, h] \;=\;
    -\tfrac{1}{2}\,r_0\,\varepsilon\, m^2
    \;+\; v\, m^6 \;-\; h\, m,
  \qquad r_0, v > 0.
\end{equation*}
The transition belongs to the $\varphi^6$ tricritical Landau universality class
(Griffiths 1970, Lawrie–Sarbach 1984), characterised by an even-parity potential
with $\mathbb{Z}_2$ symmetry $m \to -m$. This symmetry is the condition of the
theorem: its former derivation from a KO-dimension-6 real structure of the
internal triple (T-53) is retracted \status{$\times$}, since no such structure
exists on $\mathbb{C}^7$, and without it the generic codimension-3 point is the
$A_4$ swallowtail, which gives $\beta = 1/2$ (the theorem read \status{T} until
2026-09-25). UHM has exactly three physically independent
control parameters --- regeneration rate $\kappa$, dissipation rate $\gamma$, and
free-energy budget $\Delta F$ --- matching the codimension of the $\varphi^6$ tricritical
deformation $(r, u, h)$. The reduction from the 21-dimensional configuration space
to the 1-dimensional scalar order parameter $m$ is justified by the Mather Splitting
Lemma (Mather 1968): the Hessian of the effective Landau potential at the critical
point has corank 1, so $V \simeq_\mathrm{smooth} V_\mathrm{red}(m) + Q(\xi_1,\ldots,\xi_{20})$,
with $Q$ a non-degenerate quadratic in the 20 transverse modes. Mean-field exactness
follows from three independent pillars: (I) topological rigidity of the $\varphi^6$
class under smooth deformations; (II) deterministic UHM dynamics at the master-equation
(Lindblad) level, which is representation-equivalent to stochastic unraveling for
saddle-point exponents and contains no separate fluctuating field requiring renormalisation;
(III) order-parameter dimension $d_\mathrm{eff} = 21$ giving $1/N \approx 5\%$ corrections
under any stochastic reinterpretation. Note: the historical name "$A_4$-tricritical"
sometimes used in the physics literature conflates the Arnold $A_4$ swallowtail
($V_0 = x^5/5$, $\beta = 1/2$) with the $\varphi^6$ Landau tricritical ($V = rm^2/2 + vm^6$,
$\beta = 1/4$); UHM transitions belong to the latter under the $\mathbb{Z}_2$
symmetry, and to the former without it. The full derivation, including
all five exponents and verification of Rushbrooke, Widom, and Fisher scaling relations,
is given in Appendix~\ref{app:exponents}.

\paragraph{Physical meaning.}
\begin{itemize}[leftmargin=2em]
  \item $\beta = 1/4$ governs how quickly consciousness ``turns on'' above threshold.
        The value is sharper than ordinary mean-field ($\beta = 1/2$), reflecting the
        tricritical nature of the transition.
  \item $\nu = 1/2$ governs the correlation length: how far coherence extends near the
        transition. At $P \to \Pcrit$, $\xi \to \infty$ --- ``critical opalescence'' of
        consciousness.
  \item $\gamma = 1$ governs susceptibility: how sensitively the system responds to
        perturbations near threshold. At $P \to \Pcrit$, $\chi \to \infty$ --- ``critical
        slowing down.''
  \item $\alpha = 1/2$ governs the specific heat divergence: the fluctuation cost near
        threshold.
\end{itemize}

\paragraph{Comparison with known universality classes.}
The 3D Ising model has $\beta \approx 0.326$, $\nu \approx 0.630$, $\gamma \approx 1.237$.
Superfluid helium: $\beta \approx 0.347$, $\nu \approx 0.672$, $\gamma \approx 1.316$.
UHM predicts mean-field tricritical exponents $\beta = 1/4$, $\nu = 1/2$, $\gamma = 1$
--- distinct from all non-tricritical universality classes.

\begin{figure}[H]
\centering
\input{figures/critical-exponents}
\caption{\textbf{Tricritical exponent $\beta = 1/4$ vs.\ other universality
classes.} All curves are normalised so that $m/m_0 = 1$ at the reference
point $\varepsilon_0 = 0.1$, i.e.\ plotting
$y(\varepsilon) = (\varepsilon/\varepsilon_{0})^{\beta}$; this removes the
amplitude dependence $(r_0/6v)^{1/4}$ of the Landau expansion
(Appendix~\ref{app:exponents}) and isolates the exponent. The UHM
prediction $\beta = 1/4$ (thick purple) sits exactly between the two
falsification boundaries $\beta = 0.20$ and $\beta = 0.30$ (thin
dash-dotted purple), which define the acceptable UHM band. 3D Ising
$(\beta \approx 0.326)$ falls \emph{below} the lower bound for all
$\varepsilon < \varepsilon_0$, so a measured consciousness transition
with 3D Ising exponents would falsify UHM at the structural level.
Mean-field $\varphi^{4}$ $(\beta = 1/2)$ is far outside the band. The
vertical green dashed line marks $\varepsilon = 1/7$, the upper edge of
the Goldilocks window $P \in (2/7, 3/7]$ (T-124~\status{T}), beyond
which the Landau expansion is no longer physically valid.}
\label{fig:exponents}
\end{figure}

\paragraph{Falsification.}
TMS-EEG at the sleep/wake transition, $N = 50$ subjects. Fit power law $\mathrm{PCI}
\sim (P - \Pcrit)^\beta$. If $\beta \notin [0.20, 0.30]$ at $p < 0.01$ --- the
conditional theorem is falsified at the structural level (L2 falsification); a
value near $\beta = 1/2$ would single out the swallowtail branch without the
$\mathbb{Z}_2$ symmetry.

\subsection{Emergent spacetime (T-117 through T-121)}
\label{subsec:emergent-spacetime}

UHM computes an emergent four-manifold from the dynamics of $\Gam$ and obtains the
Einstein--Hilbert term from its spectral action. The chain below is \status{T} as
mathematics; its reading as physical spacetime is an interpretation \status{I}
(for $\Sigma^3$) and, for the Lorentzian sign, conditional on reflection
positivity (T-53, \status{C}). The derivation chain:

\begin{enumerate}[leftmargin=2em]
  \item \textbf{T-117 \status{T}:} Commutativity of the macroscopic algebra. In the
        thermodynamic limit ($M \to \infty$ coupled holons), macroscopic observables
        commute: $[\bar{O}_1, \bar{O}_2] \to 0$. (Proof via Goderis--Verbeure--Vets
        quantum CLT + clustering, T-39a; it holds for any local observables.)
  \item \textbf{T-118 \status{T}:} The temporal algebra $A_{\mathrm{time}} \cong C_0(\mathbb{R})$
        is the scaling limit of the reading algebras $\mathbb{C}^{N+1}$ of the
        \emph{depth register} (chronon $\Delta t \to 0$, window $\to \infty$). The
        former derivation from the summed clock $\mathbb{C}[\mathbb{Z}_{7^M}]$ of $M$
        holons is retracted \status{$\times$}: the summed clock has $6M+1$ readings and
        the fixed period $2\pi/\omega_0$; the $7^M$ ordered readings belong to the
        registers read positionally under a Feynman--Kitaev constraint.
  \item \textbf{T-119 \status{T} as mathematics (reading as physical space
        \status{I}):} For three commuting rotation charges of the holon (a maximal
        torus of $\mathrm{U}(3) = \mathrm{Stab}_{\mathrm{SO}(7)}(e_O) \cap C(L_{e_O})$;
        $\operatorname{rank}\mathrm{SO}(7) = 3$), the joint spectra of the
        fluctuations $\sqrt{M}(\bar H - \mu)$ fill $\mathbb{R}^3$, and the spatial
        algebra --- the minimal unitization --- is $C(S^3)$; the Dirac triple of
        $\Sigma^3 = S^3$ satisfies all seven of Connes' conditions. Colour-singlet
        charges give only two dimensions, so this $\Sigma^3$ has colour-charged
        coordinates (a Coleman--Mandula obstacle for the physical reading). The former
        derivation --- spectral dimension 3 from the $\{A,S,D\}$-sector via a Weyl law
        --- is retracted \status{$\times$}: a finite-dimensional $\bigotimes_m
        \mathbb{C}^3$ has no Weyl law, and the axis sector $\{A,S,D\}$ is not an
        $\mathrm{SU}(3)$ sector (row 48a).
  \item \textbf{T-120 \status{T} as mathematics:} Product spectral triple on
        $M^4 = \mathbb{R} \times S^3$; the product carries no real structure, so no
        first-order condition arises (the former KO-dimension count $4+6$ is
        retracted).
  \item \textbf{T-120b (split):} (i)~$\Sigma^3 \cong S^3$ \status{T}, now part of
        T-119 and independent of the vacuum; (ii)~constant curvature $k = +1$, de Sitter
        metric, \status{C at the vacuum symmetry}: it uses the unique sector vacuum,
        which is the hypothesis (SV) (T-64 gives no $SU(3)$-symmetric sector vacuum).
  \item \textbf{T-121 \status{T}:} Closure of Lovelock gaps 1 and 2: continuity, since
        $M^4 = \mathbb{R} \times S^3$ is smooth (T-120), and covariance, by the
        diffeomorphism invariance of the Chamseddine--Connes spectral action (the leg
        ``$G_2 \to SU(3) \to SO(3) \subset \mathrm{Diff}(M^4)$'' is retracted: no
        non-trivial homomorphism $SU(3) \to SO(3)$ exists). The spectral action yields
        the Einstein--Hilbert term; its Standard-Model part is imported from Connes'
        finite triple, not derived.
\end{enumerate}

\paragraph{New structural results.}
Fifteen results strengthen the foundations (nine new theorems + six structural upgrades):

\begin{itemize}[leftmargin=2em,nosep]
  \item \textbf{T-124b \status{T} (Independent necessity of L2 thresholds):}
        Each of $P > 2/7$, $\Phi \geq 1$, $R \geq 1/3$, $D_{\mathrm{diff}} \geq 2$
        is independently necessary for L2 consciousness. Constructive counterexamples
        demonstrate four distinct pathologies (noise-dominated, fragmented, crystallised,
        undifferentiated) when any single condition is dropped.
  \item \textbf{T-124c \status{T} (Count of nontrivial attractors; restated):}
        (1)~An isolated holon with the canonical $\vp_{\mathrm{coh}}$ has no
        stationary state besides $I/7$; (2)~with the self-registering $\vp_s$ and
        $\|H\| < h_0$ it has at least seven locally stable ones with $P > 2/7$;
        (3)~an embodied holon whose backbone rate exceeds the trace-norm Lipschitz
        constant of regeneration, $\mu > L_{\mathcal{R}}$, has exactly one, globally
        attracting at rate $\mu - L_{\mathcal{R}}$; (4)--(5)~with a constant anchor
        (the collineation anchor $\vp_J$, or any anchor with uniform diagonal) at
        $H = 0$ it has none with $P > 2/7$ below a critical $\kappa_c$ and exactly
        two above it --- a sink in $\mathcal{V}_{\mathrm{full}}$ and a saddle. The
        former statement ``at most one nontrivial fixed point in the viable set''
        is retracted \status{$\times$}: it is false by (1) and (2); its proof treated
        $\kappa(\Gam)g_V(P)$ as constants and assumed
        $\kappa_{\max} < \lambda_{\mathrm{gap}}$.
  \item \textbf{T-188 \status{C under the hypothesis T-186(a)} (Localization of the hard problem):}
        The chain $A2 \xrightarrow{\text{T-187}} J_B \xrightarrow{\text{T-185}} (\Pi \dashv \flat)
        \xrightarrow{\text{T-186(a)}} F \cong \&|_\mathcal{D}$ would reduce the hard problem of
        consciousness to ``why CPTP?'' $=$ ``why quantum mechanics?'' The chain passes
        through the hypothesis T-186(a) \status{H} (the phenomenal functor as the
        infinitesimal flat modality, asserted without a construction), so the
        localization is conditional (lowered 2026-09-25); the cohesion assumed in
        T-185 is discharged by T-185(ii$'$) \status{T}.
  \item \textbf{T-189 \status{T} (MaxEnt derivation of Bures, Char-IV):}
        Set $g_{ij} = \tfrac14 C^{\mathrm{SLD}}_{ij}$, where
        $C^{\mathrm{SLD}}_{ij} = \tfrac12\Tr(\rho\{L_i,L_j\})$ is the SLD bilinear
        form, defined without reference to any metric; Bures is the unique solution
        of this selector equation, equivalently $g_B^{-1} = 4\,C_{\mathrm{est}}$ with
        $C_{\mathrm{est}} = \mathcal{F}_{\mathrm{SLD}}^{-1}$ (Braunstein--Caves 1994:
        $\mathcal{F}_{\mathrm{SLD}} = 4g_B$). The earlier form ``$G^{ij} = C^{ij}$''
        dropped the factor 4 and inverted the tensor (corrected 2026-08-06). This
        derives $g_B$ from statistical structure without information-geometric
        choice. Char-IV is a
        \emph{physical recasting} of Char-III rather than a logically independent fourth
        witness (audit 2026-04-17): it strengthens the \emph{physical} motivation of
        T-187 without adding mathematical redundancy; T-187 retains \status{T} on Char-I
        (Petz extremality) alone.
  \item \textbf{T-190 \status{C} (Axiomatic Closure):}
        A1--A5 are derivable from (AP)\allowbreak+(PH)\allowbreak+(QG)\allowbreak+(V)\allowbreak+MaxEnt
        under three conditions: the Page--Wootters constraint $\hat C\Gam = 0$ is assumed
        (T-87, step~4), the route to A1 through T-186 is a hypothesis, and
        the enrichment of A2 is CPTP-monotone (implied by reading the enrichment as the best distinguishability over measurements, not by the holon's properties). A2 via monotonicity +
        T-187 + T-189, A3 via Theorem~S+T15, A4 from (AP) necessity, the clock register
        of A5 via T-87 steps 1--3. The earlier claim that UHM is \emph{self-grounding}
        with zero independent axioms is withdrawn (lowered 2026-09-25): the constraint
        remains an independent assumption, and the independent axiomatic content is
        A1--A4 plus the constraint of A5.
  \item \textbf{T-191 \status{T} ($\varphi$-tower convergence; restated):}
        For an embodied holon with backbone $\mu(\sigma - \Gam)$ and
        $\mu > L_{\mathcal{R}} + \kappa_{\max}$ every iterate of the tower
        ``target $a_n \to$ stationary state $a_{n+1}$'' is defined, and the tower
        converges geometrically, $q = \kappa_{\max}/(\mu - L_{\mathcal{R}}) < 1$, to
        one self-model from any anchor (T-124c(3) + Banach). The former statement
        --- convergence for every holon from any anchor --- is retracted
        \status{$\times$}: for an isolated holon the gate makes the flow bistable, so
        the iterate depends on the start. The $\varphi$-circularity is resolved by the
        closed forms of $\vp_{\mathrm{coh}}$ and $\vp_s$, not by the tower.
  \item \textbf{T-192 \status{T} ($\mathbf{Exp}^{(2)}$ is a strict 2-category):}
        The experiential 2-category satisfies all five Mac~Lane axioms (vertical/horizontal
        composition, identity 2-cells, interchange law, identity 1-cells). The lax 2-functor
        $F_2: \mathbf{DensityMat} \to \mathbf{Exp}^{(2)}$ has a valid target.
  \item \textbf{T-124d \status{T} (Threshold robustness):}
        Perturbations of order $\varepsilon$ in $\Gamma$ produce $O(\varepsilon)$
        perturbations in $P$, $\Phi$, $R$. No threshold has divergent sensitivity.
        Crossover width $\delta P \sim \varepsilon^{1/\beta} = \varepsilon^4$ (from
        $\beta = 1/4$, T-161, i.e.\ under its $\mathbb{Z}_2$ condition). The system is overwhelmingly robust for macroscopic noise
        levels.
\end{itemize}

\paragraph{Additional structural upgrades.}
\begin{itemize}[leftmargin=2em,nosep]
  \item \textbf{Octonionic correspondence upgraded $[\mathfrak{I}] \to$ \status{T}:}
        The assignment $e_k \leftrightarrow$ dimension (e.g.\ $A = e_1$, $S = e_2$, \ldots,
        $O = e_7$) is a \emph{theorem}, not an interpretation: the T15 bridge chain
        derives the octonionic structure from (AP)+(PH)+(QG)+(V) \status{T}, its
        step $\mathrm{PG}(2,2) \to \mathbb{O}$ taking the canonical orientation of
        the Fano lines --- the unique orientation class invariant under the 168
        collineations (T15-canon; the orientation input (Alt) is discharged).
        T-177~\status{T} (restated): the collineation group acts regularly on
        ordered non-collinear triples, so $O$ and the $\kappa_0$ pair $\{E,U\}$ fix
        $O$, $A$, $D$ and leave one binary choice $E \leftrightarrow U$ (with
        $L \leftrightarrow S$), which T-183 records as a convention
        \status{T}+\status{D}; the earlier uniqueness from the axis sectors of
        T-48a is retracted \status{$\times$}.
        The specific assignment is fixed up to $G_2$-gauge equivalence (T-42a~\status{T}).
  \item \textbf{Grothendieck site axioms (explicit proof):}
        The three axioms (identity, stability under pullback, transitivity) for the Bures
        site $(\mathbf{DensMat}, J_B)$ are verified via CPTP contractivity (Uhlmann 1976)
        and the essentially small presentation of compact metrizable $\Dens(\mathbb{C}^7)$
        (Johnstone, \emph{Elephant} C2.1.9--12).
  \item \textbf{7D sufficiency \status{T}:}
        The minimal 7D formalism $\Gam \in \Dens(\mathbb{C}^7)$ is sufficient for all
        consciousness observables: $P$, $R$, $\Phi$, $\CohE$, $D_{\mathrm{diff}}$, $C$;
        the 7D statements stand on their own \status{T}. They survive the round trip
        to the 42D formalism through the section--retraction $\pi \circ \iota =
        \mathrm{id}$ (T-58$'$~\status{T}). The former Morita equivalence
        $\ShInf(\mathcal{C}|_7) \simeq \ShInf(\mathcal{C}|_{42})$ (T-58) and the
        ``zero reconstruction error'' built on it are retracted \status{$\times$}:
        the equivalence fails on dimension ($\le 35$ real parameters on
        $\mathcal{D}(\ker\hat C)$ against 48 on $\mathcal{D}(\mathbb{C}^7)$).
  \item \textbf{Gap dynamics from non-associativity (Theorem 9.1 \status{T}):}
        For triples $(i,j,k)$ not on a Fano line, the octonionic associator
        $[e_i, e_j, e_k] \neq 0$ contributes to phase dynamics of coherences:
        $d\theta_{ij}/d\tau|_{\mathrm{assoc}} \propto \sum_{k} |\gamma_{ik}||\gamma_{jk}|$.
        Non-associativity is the microscopic mechanism for Gap $> 0$ (T-55).
  \item \textbf{Subjective time dilation \status{C under T-186(a)}:}
        $\delta\tau_{\mathrm{subj}}/\delta\tau_{\mathrm{phys}} \propto |\gamma_{OE}|/\gamma_{OO}$
        follows from T-87 steps 1--3 (O = clock register), T-186(a) (E = experience)
        and T-88 ($\kappa_0 \propto |\gamma_{OE}|$); since T-186(a) is now a
        hypothesis \status{H}, the result inherits its condition (it was listed as an
        upgrade to \status{T}).
  \item \textbf{The $(\mathcal{D}_\Omega, \mathcal{R})$ pair is a guiding duality
        \status{I}, not an adjunction:} $\mathcal{D}_\Omega(\Gam) = \mathrm{Hom}(\Gam,\Omega)$
        is contravariant, and the would-be Hom-set bijection fails already in
        $\mathbf{Set}$ with $\Omega = 2$ ($|\mathrm{Hom}(\mathcal{D}(1),3)| = 9 \neq 6 =
        |\mathrm{Hom}(1,\mathcal{R}(3))|$). What survives with independent
        bases: the dynamical trichotomy
        (Hamiltonian~/~dissipation~/~regeneration)~\status{T} on
        LGKS-plus-fixed-point grounds, and the $\kappa_0$
        formula~\status{T at first-order kinetics} via rapid pre-equilibrium
        (Section~\ref{sec:formalism}); the duality language organizes both.
\end{itemize}

\subsection{Foundational closure: proofs of T-188 through T-192}
\label{subsec:foundational-closure}

The five theorems T-188--T-192 form the foundational arc of UHM: they
localize the hard problem under a hypothesis (T-188), derive the metric
from first principles (T-189), derive the axioms from the characterising
properties under three conditions (T-190), bound the self-modeling tower for
embodied holons (T-191), and verify the categorical target of the
phenomenal functor (T-192). The earlier summary --- that together they
establish \textbf{zero free parameters and zero independent axioms} --- is
withdrawn: the independent axiomatic content is A1--A4 plus the
Page--Wootters constraint of A5, and the premises and free parameters
(among them the associator coupling $\kappa$ of $V_{\mathrm{Gap}}$, free by
T-331) are listed on the corpus page \emph{Premises of UHM}.

\begin{figure}[H]
\centering
\begin{tikzpicture}[
  every node/.style={font=\small},
  box/.style={draw, rounded corners, fill=white, minimum width=2.2cm,
              minimum height=0.7cm, align=center},
  arr/.style={-{Stealth[length=5pt]}, thick}
]
  \node[box, fill=axiomcolor] (AP) at (0,0) {(AP)+(PH)\\+(QG)+(V)};
  \node[box, fill=axiomcolor] (ME) at (3,0) {MaxEnt};
  \node[box] (A3) at (0,-1.5) {A3: $N\!=\!7$};
  \node[box] (A4) at (3,-1.5) {A4: $\omega_0\!>\!0$};
  \node[box] (A1) at (0,-3) {A1: $\infty$-topos};
  \node[box] (A2) at (3,-3) {A2: $J_B$};
  \node[box] (A5) at (1.5,-4.5) {A5: PW};
  \node[box, fill=theoremcolor] (T190) at (6,-2.25) {\textbf{T-190} [C]\\constraint of A5,\\T-186(a) and\\monotonicity open};
  \draw[arr] (AP) -- node[left,font=\scriptsize] {T-15} (A3);
  \draw[arr] (AP) -- node[right,font=\scriptsize] {(AP)} (A4);
  \draw[arr] (A3) -- node[left,font=\scriptsize] {T-76} (A1);
  \draw[arr] (ME) -- node[right,font=\scriptsize] {T-189} (A2);
  \draw[arr] (A1) -- node[left,font=\scriptsize] {T-87} (A5);
  \draw[arr] (A2) -- node[right,font=\scriptsize] {T-87} (A5);
  \draw[arr, dashed] (A1) -- (T190);
  \draw[arr, dashed] (A2) -- (T190);
  \draw[arr, dashed] (A3) -- (T190);
  \draw[arr, dashed] (A4) -- (T190);
  \draw[arr, dashed] (A5) -- (T190);
\end{tikzpicture}
\caption{\textbf{Axiomatic Closure (T-190) \status{C}.} A1--A4 and the clock
register of A5 are derivable from the defining properties of viable holons
(AP, PH, QG, V) plus MaxEnt, with the route to A1 through the hypothesis
T-186(a); the Page--Wootters constraint $\hat C\Gam = 0$ of A5 is not derived
and stays an independent assumption, and A2 needs, besides T-187 and T-189, the
CPTP-monotonicity of the enrichment, which the holon's properties do not give
(Lemma~M). Dashed arrows indicate that each axiom
feeds into T-190's conclusion.}
\label{fig:axiomatic-closure}
\end{figure}

\begin{theorembox}[T-188: Localization of the hard problem \status{C under the hypothesis T-186(a)}]
\label{thm:t188-localization}
The chain
\begin{equation}
  A2 \xrightarrow{\text{T-187}} J_B
  \xrightarrow{\text{T-185}} (\Pi \dashv \flat)
  \xrightarrow{\text{T-186(a)}} F \cong \&\big|_{\Dens}
  \label{eq:localization-chain}
\end{equation}
would reduce the hard problem of consciousness, within UHM, to a single
physical question: ``why experience?'' $=$ ``why CPTP?'' $=$ ``why quantum
mechanics?''. The chain passes through the hypothesis T-186(a), so the
localization is conditional (lowered 2026-09-25 from \status{T}).
\end{theorembox}

\begin{proof}
(i) T-187~\status{T}: axiom A2 uniquely determines $J_B$ (triple
characterization Char-I/II/III plus T-189 MaxEnt physical recasting).
(ii) T-185: modalities $\Pi \dashv \flat \dashv \sharp$ and
$\Re \dashv \Im \dashv \&$ exist in any differentially cohesive
$\infty$-topos \status{T}; the cohesion of the UHM setting enters through
T-185(ii$'$) \status{T} ($\mathcal{D}(\mathbb{C}^7)$ with its rank strata as
an object of $\mathrm{SynthDiff}\infty\mathrm{Grpd}$).
(iii) T-186(a)~\status{H}: the phenomenal functor $F$ is naturally
isomorphic to $\&\big|_{\Dens}$ (infinitesimal flat modality restricted
to density matrices). This step is asserted without a construction; the
earlier claim that it is ``forced by the adjunction, not stipulated'' is
withdrawn. Given A2 and T-186(a), experience \emph{is} the infinitesimal
flat modality, and the only remaining question is ``why A2?'' $=$ ``why
is the Bures metric canonical?'' $=$ ``why are CPTP maps the physical
transformations?'' $=$ ``why quantum mechanics?'' $\blacksquare$
\end{proof}

\begin{theorembox}[T-189: MaxEnt derivation of Bures (Char-IV) \status{T}]
\label{thm:t189-maxent}
Set $g_{ij} = \tfrac14\, C^{\mathrm{SLD}}_{ij}(\Gam)$, with
$C^{\mathrm{SLD}}_{ij} = \tfrac12\Tr\bigl(\Gam\{L_i, L_j\}\bigr)$ the SLD
bilinear form, defined from $\partial_i\Gam = \tfrac12(L_i\Gam + \Gam L_i)$
without reference to any metric. The Bures metric is the unique solution of
this selector equation; equivalently $g_B^{-1} = 4\,C_{\mathrm{est}}$ with
$C_{\mathrm{est}} = \mathcal{F}_{\mathrm{SLD}}^{-1}$.
\end{theorembox}

\begin{proof}
The object $\tfrac12\Tr(\Gam\{L_i,L_j\})$ \emph{is} the SLD Fisher
information $\mathcal{F}_{\mathrm{SLD}}$ (lower indices), and
$\mathcal{F}_{\mathrm{SLD}} = 4g_B$ (Braunstein--Caves~\cite{braunstein1994});
hence $g_B$ solves the selector equation. Uniqueness: distinct monotone
(Petz) means give distinct Fisher tensors, so no other monotone metric
solves it. $\blacksquare$

\noindent\emph{Correction.} An earlier version stated the selector as
$G^{ij} = C^{ij}$ (inverse metric equal to the covariance). That form is
false on two counts --- the SLD object is $\mathcal{F}_{\mathrm{SLD}}$, not its
inverse, and the factor 4 was dropped; chaining both gives $g_B^2 = \tfrac14 I$.
Uniqueness is unaffected (corrected 2026-08-06).
\end{proof}

\begin{remark}[Vanchurin correspondence]
The identification of the metric with the state's own fluctuation (SLD)
structure resonates with Vanchurin's
geometric learning theory~\cite{vanchurin2026geo}, where the metric
on the space of neural network configurations is identified with the
covariance of parameter fluctuations. In Vanchurin's framework, a
``world as a neural network''~\cite{vanchurin2020} has its geometry
determined by the statistical structure of its weights. T-189
establishes the same identification \emph{from first principles}
(MaxEnt on $\Dens(\mathbb{C}^7)$), without assuming a neural-network
substrate. The convergence of the two approaches suggests a deep
connection between geometric learning theory and UHM's categorical
framework.

The connection extends beyond the metric. In Vanchurin's later
synthesis --- the \emph{Self-Learning Universe}~\cite{vanchurin2026slu},
which derives quantum mechanics, gauge theory, and Einstein gravity as
optimality conditions of a resource-constrained distributed learning
system --- the agent replaces its three inaccessible extensive
variables (displacement $\Delta q$, fast entropy $S_x$, dof count $N$)
by their Legendre conjugates: the gauge field $A_\mu$, the temperature
$T$, and the chemical potential $\mu$. Under UHM's thermodynamic
dictionary (T-258) this triple is, channel for channel, the
three-channel perturbation basis of T-102: $A_\mu \leftrightarrow
h^{(H)}$ (the isentropic work channel), $T \leftrightarrow h^{(D)}$
(the entropy-producing heat channel; UHM's seven Fano rates refine
SLU's scalar $T$ into line-resolved temperatures whose $G_2$-symmetric
point is exactly SLU's setting), $\mu \leftrightarrow h^{(R)}$ (the
matter/feeding channel, the only entropy-importing one). The two
completeness proofs --- LGKS for T-102, the Legendre cascade over
$U(S,V,N)$ for SLU --- reach the same trichotomy from opposite ends;
the entropy--purity signatures on the UHM side are theorems (T-258),
while the identification itself is an interpretation. The two theories
are complementary along the two axes each leaves open: UHM supplies
what SLU names as its open problems (dynamical selection of the
dimension --- here $N = 7$, forced four ways --- and the non-Abelian
gauge sector --- here $G_2 = \mathrm{Aut}(\mathbb{O})$), while SLU
supplies a derivational route to what UHM postulates (the
density-matrix substrate as the coarse-grained limit of a classical
learning system). Two later corpus results sharpen the dictionary:
the regeneration channel is the \emph{unique optimal} learning flow
(steepest descent along the mixture geodesic of the Kubo--Mori
metric, the unique dually flat monotone metric --- Grasselli--Streater;
corpus T-263), and SLU's ``gravity is the efficiency of learning''
acquires sign and sectors through the information--gravity reciprocity
$G_N^{(ST)}\langle\mathrm{QFI}(\theta_{\mu\nu})\rangle_{\mathrm{ST}} =
8\pi N\mu^2$ with $\Lambda \propto \mathcal{G}_O^2$ (corpus T-264,
T-254): learnable phases gravitate weakly, and the unlearnable clock
phase is the cosmological constant.
\end{remark}

\paragraph{The five characterising properties.}
T-190 below uses five properties that define a viable holon. For the
reader's reference, they are collected here:

\begin{defbox}[Characterising properties of a viable holon \status{D}]
\label{def:five-properties}
\begin{enumerate}[leftmargin=2em,nosep]
  \item \textbf{(AP) Autopoiesis:} the system maintains $P > \Pcrit$
        against dissipation via internal regeneration $\Regen$.
  \item \textbf{(PH) Phenomenal identification:} the internal aspect
        of $\Gam$ \emph{is} experience (two-aspect monism,
        Def.~\ref{def:two-aspect}).
  \item \textbf{(QG) Quantum-gravitational consistency:} the system's
        dynamics are compatible with the CPTP structure of quantum
        mechanics (the evolution preserves $\Gam \in \Dens(\mathbb{C}^7)$).
  \item \textbf{(V) Viability:} the system has a non-trivial stationary
        state $\rhostar \neq I/7$ (it is alive, not thermally dead).
  \item \textbf{(MaxEnt) Maximum entropy:} the system's fluctuation
        structure is Jaynes-maximal --- the metric equals one quarter of
        the SLD form, $g_{ij} = \tfrac14 C^{\mathrm{SLD}}_{ij}$ (T-189~\status{T}). This is a physical
        recasting of SLD-Cram\'er--Rao saturation (Char-III), not an
        additional logically independent constraint.
\end{enumerate}
These are not axioms but \textbf{definitions}: they specify what it
means for a physical system to be a viable self-sustaining holon. T-190
derives A1--A4 and the clock register of A5 from (AP)--(MaxEnt) under its
three conditions. A complementary
\emph{inverse} closure T-181~\status{T} (Bimodule construction) derives
(AP), (PH), (QG), (V) themselves from the axioms A1--A4: (QG) from A1,
(AP) from A1, (PH) from A1+A3, (V) from A2+A3. The pair (T-181, T-190)
relates ``axioms'' and ``characterising properties'' in both directions,
but it is not a full equivalence: the Page--Wootters constraint of A5,
the hypothesis T-186(a) and the monotonicity of the enrichment of A2 stay
outside the circle, so UHM is not
self-grounded (the earlier claim is withdrawn with the lowering of T-190).
\end{defbox}

\begin{theorembox}[T-190: Axiomatic Closure \status{C}]
\label{thm:t190-closure}
Under the constraint assumption of T-87 (step~4), the hypothesis
T-186(a) and the CPTP-monotonicity of the enrichment of A2 (Lemma~M:
implied by the operational reading of the enrichment, not by
(AP)+(PH)+(QG)+(V)+MaxEnt; named 2026-09-26), A1--A5 are derivable from (AP)+(PH)+(QG)+(V)+MaxEnt
(Figure~\ref{fig:axiomatic-closure}). The converse direction
(A1--A4 $\Rightarrow$ (AP)+(PH)+(QG)+(V)) is T-181~\status{T}. The earlier
statement ``UHM has zero independent axioms'' is withdrawn (lowered
2026-09-25): the constraint $\hat C\Gam = 0$ remains an independent
assumption, and its independence of A1--A4 is \status{T}.
\end{theorembox}

\begin{proof}
\emph{A3} ($N = 7$): Theorem~S (functional minimality) derives the
\emph{number} $N = 7$ from (AP)+(PH)+(QG); the bridge chain
T-15~\status{T given Theorem S} then forces the \emph{structure} of the
seven-dimensional system --- $\mathrm{BIBD}(7,3,1) \to \mathrm{PG}(2,2) \to
\mathbb{O}$ --- and $\dim\mathrm{Im}(\mathbb{O}) = 7$ (Hurwitz--Adams) closes
the consistency loop on the same number.
\emph{A4} ($\omega_0 > 0$): $\omega_0 = 0 \Rightarrow H_{\mathrm{eff}} = 0
\Rightarrow$ no dynamics $\Rightarrow$ (AP) violated. Hence $\omega_0 > 0$
is necessary for autopoiesis.
\emph{A1} ($\infty$-topos): T-76~\status{T} at site level: the Bures
topology on $\Dens(\mathbb{C}^7)$ generates a Grothendieck site whose
sheaves form an $\infty$-topos (Lurie HTT~6.2.2.7); the route from the
operational basis to A1 goes through T-186(a), a hypothesis \status{H}.
\emph{A2} ($J_B$) \status{C} (monotonicity): the topology is forced (every
continuous distance on the compact $\Dens(\mathbb{C}^7)$ induces the standard
one); \emph{within the CPTP-monotone metrics} T-187~\status{T} +
T-189~\status{T}: triple mathematical characterization (Char-I Petz extremality
+ Char-II Uhlmann + Char-III SLD-Fisher) plus MaxEnt physical recasting
(Char-IV) uniquely selects Bures. The monotonicity is not derived: the
Hilbert--Schmidt distance, in which (V) is written, satisfies every property
and MaxEnt, yet a partial-trace channel on $\mathbb{C}^7$ moves a full-rank
pair $\sqrt2 = 1.41421$ times farther apart in it while their Bures angle keeps
the ratio $1.00000$. Under the operational reading (O),
$d(\rho,\sigma) = \sup_M \delta\bigl(p_M(\rho), p_M(\sigma)\bigr)$ over POVMs
$M$, monotonicity is a theorem (the adjoint of a channel maps POVMs to POVMs),
and (O) with the Bhattacharyya angle gives the Bures angle
directly~\cite{fuchs1995} (Lemma~M; this step read \status{T} until
2026-09-26).
\emph{A5} (Page--Wootters): T-87, steps 1--3~\status{T}: the Wedderburn
decomposition of $A_{\mathrm{int}}$ isolates the clock summand and the tensor
factor $\mathcal{H}_O = \mathbb{C}[\mathbb{Z}_7]$, so the clock register follows
from A1--A4. Step~4, the constraint $\hat C\Gam = 0$, is \status{C} under the
support condition of Property~2: stationarity of a mixed state gives only
$[\hat C, \Gam] = 0$, and the earlier ``A5 is a consequence of A1--A4'' is
withdrawn. $\blacksquare$
\end{proof}

\begin{theorembox}[T-191: $\vp$-tower convergence, restated \status{T}]
\label{thm:t191-convergence}
Let the holon be embodied, with backbone $\mu(\sigma - \Gam)$ and
$\mu > L_{\mathcal{R}} + \kappa_{\max}$, where $L_{\mathcal{R}}$ is the
trace-norm Lipschitz constant of regeneration, uniform in the target. Then
every iterate of the tower ``target $a_n \to$ stationary state $a_{n+1}$''
is defined, and the tower converges geometrically to one self-model from
any anchor:
\begin{equation}
  \|a_{n} - a^*\|_1 \leq q^n\,\|a_0 - a^*\|_1,
  \quad q = \frac{\kappa_{\max}}{\mu - L_{\mathcal{R}}} < 1.
  \label{eq:phi-tower-convergence}
\end{equation}
\end{theorembox}

\begin{proof}
Under backbone dominance each target has exactly one stationary state,
globally attracting at rate $\mu - L_{\mathcal{R}}$ (T-124c(3)); a
variation-of-constants estimate bounds the dependence of that state on
the target by $q$, and Banach's fixed-point theorem gives the geometric
convergence. $\blacksquare$

\noindent\emph{Erratum \status{$\times$}.} The former statement ---
convergence for every holon from any anchor with
$q = \kappa_{\max}/(\lambda_{\mathrm{gap}} + \kappa_{\min}) < 1$ by T-39a +
T-96 --- is retracted: for an isolated holon the gate $g_V$ makes the flow
bistable, so the iterate $\lim_\tau \exp(\tau\Lind^{(n)})$ depends on the
start (from $I/7$ it stays dead, from $|0\rangle$ it lives), and T-96 gives
no bound $\kappa_{\max} < \lambda_{\mathrm{gap}}$. The $\vp$-circularity is
resolved by the closed forms of $\vp_{\mathrm{coh}}$ and $\vp_s$, not by
the tower.
\end{proof}

\begin{theorembox}[T-192: $\mathbf{Exp}^{(2)}$ is a strict 2-category \status{T}]
\label{thm:t192-2cat}
The experiential 2-category $\mathbf{Exp}^{(2)}$ satisfies all five
Mac\,Lane axioms: vertical/horizontal composition, identity 2-cells,
interchange law, and identity 1-cells.
\end{theorembox}

\begin{proof}
(i)~Vertical composition: for 2-cells $\alpha: \Phi \Rightarrow \Psi$
and $\beta: \Psi \Rightarrow \chi$, pointwise $(\beta \circ_v \alpha)_\rho
:= \beta_\rho \circ \alpha_\rho$ is natural (Mac\,Lane CWM IV.2).
(ii)~Horizontal composition: Godement product
$(\beta \circ_h \alpha)_\rho := \beta_{\Phi_2(\rho)} \circ \Psi_1(\alpha_\rho)$.
(iii)~Identity 2-cells: $(\mathrm{id}_\Phi)_\rho = \mathrm{id}_{\Phi(\rho)}$.
(iv)~Interchange: $(\beta_2 \circ_v \beta_1) \circ_h (\alpha_2 \circ_v \alpha_1)
= (\beta_2 \circ_h \alpha_2) \circ_v (\beta_1 \circ_h \alpha_1)$
by the Eckmann--Hilton argument.
(v)~Identity 1-cells: $F(\mathrm{id}_\rho) = \mathrm{id}_{F(\rho)}$
(functoriality of $F$). All hold with equalities (strict), so
$\mathbf{Exp}^{(2)}$ is strict. The lax 2-functor $F_2$ has a valid
target. $\blacksquare$
\end{proof}

\subsection{SYNARC-derived theorems (T-193--T-209)}
\label{subsec:synarc-theorems}

The SYNARC companion specification (Appendices G, H, I) establishes a
layered chain of meta-theorems culminating in AGI- and ASI-sufficiency
and operational deployability. Full proofs are in the extended
documentation; their statements are collected here as dependencies for
the fundamental closures of Section~\ref{subsec:fundamental-closures}.

\begin{itemize}[leftmargin=2em,nosep]
  \item \textbf{T-193~\status{T} (Yoneda universal representability):}
        every CPTP-implementable task has a representable sheaf in
        $\ShInf(\Dens(\mathbb{C}^7), J_{\mathrm{Bures}})$. Upgraded to a
        computable, Kolmogorov-free formulation in T-213 below.
  \item \textbf{T-194~\status{T} (Cram\'er--Rao saturation):}
        Bures-gradient learning attains the quantum Cram\'er--Rao bound
        up to a constant factor ($C_2 \leq 16$; the earlier $C_2 \leq 4$ did not
        follow from the stated ingredients, corrected 2026-08-06).
  \item \textbf{T-195~\status{T} (L-III $\Phi$-monotonicity):}
        refinements of the epistemic topology are non-decreasing in
        $\Phi$; strengthened to strict monotonicity on viable states in
        T-210 below.
  \item \textbf{T-196~\status{T} (Goldilocks sustainability):}
        closed sensorimotor loops preserve $P \in (2/7, 3/7]$ under
        bounded perturbations.
  \item \textbf{T-197~\status{T}+\status{D} (AGI-sufficiency, S-11):}
        any realisation of the SYNARC architecture (definition~\status{D})
        satisfies the UHM-AGI predicate (A1)--(A7); the derivation
        chain~\status{T} is non-tautological because each architectural
        component is forced by a distinct load-bearing theorem. Caveat:
        A7 (recursive self-improvement) is \emph{weakly} monotone by
        T-195; strict improvement requires T-210 on interior states.
  \item \textbf{T-198~\status{T} (G\"odelian creativity):}
        strictly monotone ordinal functors $\mathfrak{A}_\bullet:
        \mathrm{On} \to \mathbf{Cat}_\infty$ with $G_2$-preserving
        inclusions admit, at each step, a representable sheaf
        non-equivalent to any in the previous layer.
  \item \textbf{T-199--T-202 (Value structure, L-IV
        modification, KS contextuality, meaning as $G_2$-orbit):}
        SYNARC Appendix H, Theorems H.2--H.5; T-199, T-200~\status{T},
        T-202~\status{T}+\status{I}. T-201 in its original form ---
        Fano-line projectors admit no joint distribution --- is retracted
        \status{$\times$}: those projectors are diagonal in the pointer basis
        and commute, and $p_i = \gamma_{ii}$ reproduces every context. The
        corrected T-201$'$~\status{T}: the 63 rays of the $E_7$ configuration
        in $\mathbb{R}^7$, built from the Fano plane and the Hamming code, form
        135 orthonormal bases admitting no Kochen--Specker assignment.
  \item \textbf{T-203~\status{T}+\status{I} (Qualia as Gap spectral
        eigenvectors):} the mathematical core (eigenvector class of
        $\hat{\mathcal G}|_E$, $G_2$-covariant and Gap-faithful) is
        \status{T}; the identification
        $\mathrm{Qualia}(\Gam) := \text{eigenvector class}$ is
        \status{I} (ontological postulate bridging structure to
        phenomenology), mirroring the T-38a stratification.
  \item \textbf{T-204~\status{T} (Pareto-optimal bounded rationality):}
        graceful degradation of UHM-AGI under the resource budget
        $(C, M, \varepsilon)$; at $d_{\mathrm{eff}} = 2$ the system drops
        to the minimal-consciousness floor $\Dmin = 2$.
  \item \textbf{T-205~\status{C}+\status{D} (Ordinal mentalisation via
        fractal-holon tower):} the claim
        $\omega^\omega$-depth holds for $\iota_{\max}$ identity
        convention with abstracted resource constraints, contingent on
        (i)~bounded Landauer cost C22, (ii)~cocompleteness of the
        filtered colimit along an $\omega^\omega$-chain in
        $\ShInf(\mathcal C)$, (iii)~the interpretive choice of a
        single-agent identity predicate. Finite truncation is
        unconditionally \status{T}; the identity choice is made explicit
        in T-215 below.
  \item \textbf{T-206--T-209 (Operational protocols, S-13):}
        qualia tomography, inverse value-alignment, constructive
        existence of $G_2$-invariant value sets (T-206--T-208~\status{T};
        in T-208 the integration-based family is invariant only under the
        finite frame group $\Gamma_{\mathrm{oct}}$, since $\Phi$ is
        frame-pinned), and the operational closure meta-theorem
        (T-209~\status{T}+\status{D}; SYNARC Appendix I).
\end{itemize}

\subsection{Fundamental closures: T-210--T-216}
\label{subsec:fundamental-closures}

The following seven theorems address mathematical and categorical gaps
of the UHM framework; of the fourteen closures T-210--T-223, eleven stand
as \status{T} (T-211 and T-212 in corrected forms), T-221 is stratified
into \status{T} and \status{I} parts, T-216 is \status{C} and T-219 is
\status{H}. The earlier summary ``no open mathematical or categorical gaps
remain'' is retracted \status{$\times$}: the \status{C} and \status{H} rows
are open conditions. Full proofs are collected in the
companion document \emph{Fundamental Closures}; the statements are
reproduced here for self-containment.

\begin{theorembox}[T-210: Strict $\Phi$-monotonicity on interior
states \status{T}]
\label{thm:t210-strict-phi}
Let $J \subsetneq J'$ be a proper refinement of the Bures-compatible
epistemic topology. For any state $\Gam$ in the interior stratum
$\mathcal D_7$ (full rank, all $|\gamma_{ij}| > 0$),
\[
  \Phi(\Gam \mid J') - \Phi(\Gam \mid J) \;\geq\;
  \frac{\min_{(i,j) \in J' \setminus J}|\gamma_{ij}|^2}
       {\sum_k \gamma_{kk}^2} \;>\; 0.
\]
\end{theorembox}

\begin{proof}
$\Phi(\Gam \mid J) = \bigl(\sum_{(i,j) \in \mathrm{supp}(J) \cap
\mathrm{Off}} |\gamma_{ij}|^2\bigr) / N_\Gam$ with $N_\Gam =
\sum_k \gamma_{kk}^2$. On $\mathcal D_7$, every off-diagonal contributes
strictly positively; proper refinement adds at least one pair
$(i^*, j^*)$ with $|\gamma_{i^*j^*}|^2 > 0$. Since $\mathrm{supp}(J)
\subsetneq \mathrm{supp}(J')$ and $\Gam$ has full rank on $\mathcal D_7$,
the interior hypothesis is non-vacuous. $\blacksquare$
\end{proof}

Consequence: T-195 is upgraded from \emph{weak} to \emph{strict}
monotonicity for states in the interior stratum $\mathcal{D}_7$ (full
rank, all $|\gamma_{ij}| > 0$) whenever the agent crosses a proper
L-III refinement $J \subsetneq J'$. The A7 clause of T-197 thus implies
strict self-improvement for every viable SYNARC agent \emph{in the
interior stratum, under L-III refinement events crossing the
obstruction threshold $\omega(J_{\mathrm{ep}}) > \omega_{\mathrm{th}}$};
continuous strict improvement outside these conditions remains
\status{C} (pending T-197 refinement).

\begin{theorembox}[T-211: PhysTheory is an $(\infty,1)$-category \status{T}]
\label{thm:t211-phystheory}
$\mathbf{PhysTheory}$ (physical theories with CPTP dynamics, T-174) is
the cartesian unstraightening of $E \mapsto \mathrm{Fun}(B\mathbb{R},
\mathrm{Alg}(E))$ over Lurie's $\mathbf{Topoi}_\infty$ (HTT~\S3.2), a
Grothendieck construction. It is a quasicategory, so associativity up to
coherent homotopy, the pentagon, the interchange law and all higher
simplicial identities hold; the forgetful functor to
$\mathbf{Topoi}_\infty$ is not faithful.
\end{theorembox}

\begin{proof}
Straightening/unstraightening (HTT Thm.~3.2.0.1) turns the functor
$E \mapsto \mathrm{Fun}(B\mathbb{R}, \mathrm{Alg}(E))$ into a cartesian
fibration over $\mathbf{Topoi}_\infty$, whose total space is a
quasicategory; the mapping space over a geometric morphism $f$ is
$\mathrm{Map}_{\mathrm{Dyn}(E_1)}(x_1, f^*x_2)$, i.e.\ the triples
$(f^*, \alpha, \beta)$ of T-174. Complex conjugation of $(\mathbb{C},\cdot)$
over $\mathrm{id}_{\mathcal S}$ shows the forgetful functor is not faithful.
$\blacksquare$

\noindent\emph{Erratum \status{$\times$}.} The earlier statement ---
$\mathbf{PhysTheory}$ is a \emph{full} $(\infty,1)$-subcategory of
$\mathbf{Topoi}_\infty$, fully faithful via T-173, with coherences
inherited through HTT~5.2.7 --- is retracted; its proof also invoked
T-178 (retracted as a derivation) and the reconstruction of T-119, which
the corrected statement does not need. (Status history: \status{T}, then
\status{C at T-119}, then \status{T} in the corrected form, 2026-09-25.)
\end{proof}

\begin{theorembox}[T-212$'$: U-projection --- the $G_2$-twirl \status{T}]
\label{thm:t212-rh}
For $G_2$ acting on $\mathbb{C}^7$ by its seven-dimensional
representation, $\int_{G_2} gXg^\dagger\,dg = \tfrac17\Tr(X)\,\mathbf 1$:
the commutant is $\mathbb{C}\mathbf 1$ (Schur), and this is the unique
trace-preserving map onto $\mathbb{C}\mathbf 1$, idempotent only with the
normalisation $\tfrac17$. Its reading as the U dimension (Unity) is an
interpretation \status{I}.
\end{theorembox}

\noindent\emph{Erratum \status{$\times$}.} T-212 earlier identified this
map with the rheonomy modality Rh, ``the seventh canonical modality of
T-185'', right adjoint to the bosonic-grade forgetful functor in a
super-cohesive extension (citing DCCT~\cite{schreiber2013} \S3.10). The
identification is retracted: Rh of solid cohesion preserves global points,
whereas the formula sends every state to $I/7$; Rh is not among the seven
modalities of differential cohesion (T-185(iii)); it does not occur in DCCT
v1, whose \S3.10 is ``Structures in a differentially cohesive
$\infty$-topos''; and idempotence fails with the unnormalised trace. The
former bijection of the seven modalities with the seven dimensions
($\mathrm{Id} \leftrightarrow O$, $\Pi \leftrightarrow A$,
$\flat \leftrightarrow S$, $\Im \leftrightarrow D$,
$\sharp \leftrightarrow L$, $\& \leftrightarrow E$,
$\mathrm{Rh} \leftrightarrow U$) is an interpretation \status{I} at most,
with Rh replaced by Red ($\Re$) (T-185(iii)).

\begin{theorembox}[T-213: Yoneda representability via Bures
description length \status{T}]
\label{thm:t213-bures-yoneda}
For any CPTP-implementable $f: \mathrm{Obs} \to \mathrm{Act}$, define
$D_B(f) := \min_{\rho_f \text{ implementing } f}
|\mathrm{Kraus}(\rho_f)| \cdot \log_2 7 \in \mathbb N \cdot \log_2 7$.
Then the representable sheaf $F_f$ satisfies
\[
  \|F_f\|_B \;\leq\; C_1 \cdot D_B(f) \cdot \log(1/\varepsilon),
  \qquad C_1 = \omega_0^{-1}\log 7,
\]
with the universal bound $D_B(f) \leq 49\log_2 7 \approx 138$ bits by
Stinespring dilation.
\end{theorembox}

This upgrades T-193 to a \emph{computable}, Kolmogorov-free form: the
original bound $\|F_f\|_B \leq C_1 K(f) \log(1/\varepsilon)$ depended on
the uncomputable Kolmogorov complexity $K(f)$; the Bures description
length $D_B(f)$ is computable for every $f$ and universally bounded.

\begin{theorembox}[T-214: Hard-problem meta-theorem \status{T}]
\label{thm:t214-hard-problem}
Let $W: \Dens(\mathbb C^7) \to \mathrm{Mind}$ be a putative bridge
functor assigning experiential content to coherence states. Then:
(i)~$W$ cannot be expressed as a morphism internal to
$\mathrm{Th}_{\mathrm{UHM}}$ without violating Lawvere's fixed-point
theorem and T-55;
(ii)~the identifications ``E-sector = interiority'' (T-38a) and
``qualia = eigenvectors'' (T-203) are \emph{necessarily} external
postulates \status{P}/\status{I};
(iii)~this is a \emph{positive} structural result, not a lacuna.
\end{theorembox}

\begin{proof}[Proof sketch]
If $\tilde W \in \Omega^{\mathcal D}$ is internal and point-surjective,
the self-referential predicate
$\Phi(\Gam) := \neg \exists \Gam': \tilde W(\Gam') = \tilde W(\Gam)$
admits a Lawvere fixed point contradicting itself. Hence no internal
$\tilde W$ can simultaneously be faithful to external content and
surjective. The bridge must be external. $\blacksquare$
\end{proof}

Combined with T-188 (WHY-localisation to ``why CPTP?'', conditional on
the hypothesis T-186(a)) and T-203 (WHAT-structure up to $G_2$-gauge),
T-214 shows that the residual \status{P}/\status{I} status
of phenomenal identification is structurally inevitable, not
a remediable weakness. Further internal progress is impossible by
Lawvere incompleteness, and this is positively proven.

\begin{theorembox}[T-215: Cross-layer identity convention
\status{T}+\status{D}]
\label{thm:t215-cross-layer}
For a fractal holon tower $\mathcal T = (A_0, A_1, \ldots)$ with
$A_{n+1}$ extending $A_n$ via \texttt{spawn\_child}, the predicate
``$\mathcal T$ is a single agent'' is conventionally determined by a
choice $\iota \in \{\iota_{\min}, \iota_{\max}\}$: $\iota_{\min}$
(society, per-agent $\mathrm{SAD} \leq 3$) or $\iota_{\max}$
(composite, unbounded cross-layer depth subject to C22 Landauer
cost and T-204 bounded rationality). Both are consistent with
$\Omega^7$; neither is derivable from the axioms alone.
\end{theorembox}

This explains the apparent tension between SAD\textsubscript{MAX}=3
(per-agent, T-142~\status{T}) and the $\omega^\omega$-depth of T-205:
the latter holds only for the $\iota_{\max}$ convention with
resource abstraction.

\begin{theorembox}[T-216: Closed-form analytical
$\varepsilon_{\mathrm{eff}}$ \status{C at (SV)}]
\label{thm:t216-eps-eff}
Under the hypothesis (SV) of sector values of the Gap vacuum, the
effective sectoral parameter in the Yukawa hierarchy has the symbolic
form
\[
  \varepsilon_{\mathrm{eff}} \;\propto\;
  \frac{N_{33}^{\mathrm{Fano}}}
       {9\, |\bar\gamma|\,\bigl(1 + r_4\,\Sigma_0/2\bigr)},
\]
with $N_{33}^{\mathrm{Fano}}$ the number of non-$O$ Fano lines meeting
the $\mathbf 3$-sector in exactly two points, $|\bar\gamma|$ the
sectoral average of off-diagonal coherences at the vacuum $\theta^*$,
$r_4 = V_4/V_2|_{\theta^*}$, and $\Sigma_0 = \sum_i \theta_i^{*2}$.
The value $\varepsilon_{\mathrm{eff}} \approx 0.059$ is the
phenomenological sectoral value, \status{C}, pending a corrected
symbolic evaluation.
\end{theorembox}

Derivation: symbolic minimisation of $V_{\mathrm{Gap}}$ (T-74~\status{T})
in the sectoral variables $\boldsymbol\varepsilon$, with the Fano
selection rule (T-43d~\status{T}) and the Schur-unique $G_2$-invariant
cubic term (T-50~\status{T}) supplying the combinatorial content.
\emph{Corrections.} (1)~The earlier ``$N_{33}^{\mathrm{Fano}} = 1$, the
unique Fano line $\{L,E,U\}$ within $\bar{\mathbf 3}$'' is false:
$\{L,E,U\}$ is not a line, and no line lies wholly within
$\bar{\mathbf 3}$. (2)~The printed $0.059$ does not follow from the stated
$|\bar\gamma| \approx 0.15$ (the ratio is $O(1)$). (3)~The status was
\status{T at T-64}; the sectoral reduction rested on the axis-labelled
decomposition T-48a, retracted, and T-64 as corrected produces no sector
pattern, so the structure rests on (SV) (corrected 2026-09-25).

\subsection{Extended closures: T-217--T-223 + clarifications}
\label{subsec:extended-closures}

A follow-up audit (2026-04-17) identified three further open
mathematical questions and three implicit assumptions; T-217 and T-220
are \status{T}, T-218 is \status{T} for the Kan property with its
3-coskeletal step \status{C}, and T-219 is a hypothesis \status{H}
(corrected 2026-09-25). A further theorem (T-221, 2026-04 v1.5; corrected
2026-09-25) locates UHM with respect to the recent no-go results of List
(2025)~\citep{list2025quadrilemma} and DeBrota--List
(2026)~\citep{debrota2026heptalemma,debrotalist2026preprint}. The
resource-theoretic theorem T-222 (restated 2026-09-26; the former
``MRQT-completeness'' version of 2026-04-18 is retracted) establishes
that the viability window selects no single resource optimum: every
viable state is strictly dominated on every R\'enyi free energy by
partial depolarisation, the Pareto set lies on the sphere $P = 2/7$,
and $F_1$ and $F_\infty$ are minimised by different spectra. The
Putnam-triviality closure T-223 (2026-04-18 v1.7)
responds to Lerchner (2026) ``The Abstraction Fallacy'' by
foreclosing Putnam-style alphabetization indeterminacy at the
intrinsic categorical layer $L2 = [\Gamma]_{G_2}$.

\begin{theorembox}[{T-217: L3 tricategorical coherence \status{T}}]
\label{thm:t217-l3-tricat}
The third-level interiority category $\mathbf{Exp}^{(3)} :=
\tau_{\leq 3}(\mathbf{Exp}_\infty)$ is a coherent tricategory in the
Gordon--Power--Street sense. The $K = 3 + 1 = 4$ decomposition is
derived: three inherited 2-cells (LGKS triadic: Aut, $\mathcal D$,
$\mathcal R$) plus one new 3-cell modification
$\eta : \varphi^{(2)} \Rightarrow \varphi \circ \varphi$.
\end{theorembox}

\begin{proof}[Proof sketch]
(1)~By T-91~\status{T}, $\mathbf{Exp}_\infty = \mathrm{Sing}(\mathcal E)$
is a Kan complex (Milnor 1957). (2)~$\tau_{\leq 3}$ yields a 3-truncated
Kan complex, i.e.~a 3-type (Lurie HTT 5.5.6.18). (3)~3-types are
equivalent to coherent tricategories with invertible cells
(Baez--Dolan hypothesis, proved for $n \leq 3$ by Hirschowitz--Simpson
2001, Leinster 2002; Gordon--Power--Street 1995 coherence).
(4)~The cellular count $K = 3 + 1$ follows from T-57~\status{T} (LGKS
three 2-cells) and the strict-2-category enrichment T-192~\status{T}
(unique non-trivial 3-cell is the coherence modification $\eta$).
Pentagon-of-pentagons holds automatically for $\tau_{\leq 3}$ of a Kan
complex (Lurie HTT 5.2.7). $\blacksquare$
\end{proof}

\textbf{Upgrade of T-67.} The ``$3+1$ decomposition'' in the Bayesian
derivation of $R^{(2)} \geq 1/4$ (L3 threshold) is now
\emph{derived from tricategorical first principles}, not heuristic.
T-67 is thus [T], with all justification in T-217 above.

\begin{theorembox}[{T-218: SYNARC cognitive complex is a Kan complex
\status{T}}]
\label{thm:t218-cog-kan}
The SYNARC cognitive simplicial set
$\mathrm{Cog} := \mathrm{Sing}(B_\bullet \mathcal C_{\mathrm{FKraus}})$,
defined as the singular complex of the classifying space of the
Fano-Kraus category, is a Kan complex: every horn $\Lambda^n_k \to
\mathrm{Cog}$ (inner or outer) admits a filler $\Delta^n \to \mathrm{Cog}$.
Its 3-coskeletal truncation $\tau_{\leq 3}\mathrm{Cog} \simeq \mathrm{Cog}$
matches $\mathrm{SAD}_{\max} = 3$ at the categorical level --- \status{C},
on the SYNARC-viable subset only (the step bridging simplicial and
Bures-metric viability is not a simplicial identity; off that subset
$\tau_{\leq 3}$ is the ordinary truncation and no equivalence).
\end{theorembox}

\begin{proof}[Proof sketch]
Milnor 1957: for any topological space $X$, $\mathrm{Sing}(X)$ is a Kan
complex. Apply to $X = B_\bullet \mathcal C_{\mathrm{FKraus}}$ (the
classifying space of the Fano-Kraus category; a CW-complex by
Segal 1968). An explicit horn-filler algorithm is given via radial
pullback $r_k : \Delta^n_{\mathrm{top}} \to \Lambda^n_k$ with complexity
$O(n \cdot \dim \mathcal D)$. The 3-coskeletal truncation is a
3-truncated Kan complex (Lurie HTT 5.5.6.21), matching T-142 SAD
$_{\max}$. $\blacksquare$
\end{proof}

\textbf{Replaces the nerve construction.} Previously SYNARC's Cog was
identified as the nerve $N(\mathcal C_{\mathrm{FKraus}})$ of a 1-category,
which satisfies the inner Kan condition but not outer horn filling in
general. The singular-classifying-space construction above closes this
gap with a textbook-level application of algebraic topology.

\begin{theorembox}[{T-219: $\Lambda$ SUSY-suppression via sector product
\status{H}}]
\label{thm:t219-lambda-susy}
\emph{Hypothesis} (corrected 2026-09-25 from \status{T at T-64}). In UHM's
$\mathcal N = 1$ spectral action on $M^4 \times A_{\mathrm{int}}$
(T-65), the residual cosmological constant from SUSY-broken
loops would be suppressed by
\[
  \Lambda_{\mathrm{SUSY}} \;\sim\; \varepsilon^{12}\, M_P^4 \;=\;
  \varepsilon^{4 \cdot k_{\mathrm{sec}}}\, M_P^4, \qquad k_{\mathrm{sec}} = 3,
\]
where $k_{\mathrm{sec}} = 3$ is the number of sectors
($\mathbf 1_O \oplus \mathbf 3 \oplus \bar{\mathbf 3}$)
and the factor $4$ per sector is the standard one-loop supertrace
$\mathrm{STr}(M_k^4) \sim (\varepsilon M_P)^4$ of SUSY-broken
spectrum-splitting (Martin 2010, \emph{SUSY Primer} \S7.2).
\end{theorembox}

\noindent Why a hypothesis: the sectors were the axis triples of T-48a,
retracted \status{$\times$} in its axis-labelled form; their breaking scales
rest on (SA) \status{H} and (FE); and the proof's own one-loop sum
$\sim 3\varepsilon^4 M_P^4$ exceeds $\varepsilon^{12} M_P^4$ unless the
lower orders cancel, which is not shown.

\textbf{Error correction.} The previous formulation claimed
12-order suppression via decomposition of the $G_2$ adjoint
$\mathbf{14} \to \mathbf{7}_{\mathrm{light}} \oplus
\mathbf{7}_{\mathrm{heavy}}$. This is \emph{mathematically invalid}:
the $G_2$ adjoint representation $\mathbf{14}$ is irreducible under
$G_2$, and no $7+7$ decomposition exists. T-219 replaces this with a
sector-product mechanism: the three sectors of the UHM state
space would each contribute one factor of $\varepsilon^4$ in the leading
three-loop correction, giving $\varepsilon^{12}$ via
$G_2$-invariant Fano coupling (T-43d~\status{T}) instead of invoking a
reducibility that does not hold --- a hypothesis for the reasons above,
not a rigorous derivation.

\paragraph{Three clarifications (implicit assumptions made explicit).}
\begin{itemize}[leftmargin=2em,nosep]
  \item \textbf{A4 simple spectrum.} $H_{\mathrm{eff}}$ has seven
        distinct eigenvalues (generic condition; required for
        $\mathbb Z_7$ temporal modality and Berry-phase calculations on
        non-degenerate strata).
  \item \textbf{$f_0$ canonical formula.} The spectral zeta
        $\zeta'_{H_{\mathrm{Gap}}}(0) = -\log \det H_{\mathrm{Gap}}$
        is well-defined on the finite-dimensional $(S^1)^{21}/G_2$
        phase space without regularisation ambiguity (elementary
        rational expression in eigenvalues).
  \item \textbf{Bures boundary.} Rank-deficient boundary strata of
        $\mathcal D(\mathbb C^7)$ are handled via the stratified site
        (Ayala--Francis--Rozenblyum 2017). The viability condition
        $P > \Pcrit$ together with $D_{\min} \geq 2$ (an independent L2
        threshold \status{D}, T-151)
        restricts attention to the interior stratum $\mathcal D_7$
        where the Bures metric is non-degenerate. Boundary handling is
        needed only for pathological-state analysis.
\end{itemize}

\begin{theorembox}[{T-220: No-reduction theorem for $F_4$-UHM $\to$
$G_2$-UHM \status{T, negative}}]
\label{thm:t220-no-reduction}
Let $\mathbf{C}_{F_4}$ denote the hypothetical base category of
$F_4$-UHM --- states on the exceptional Jordan algebra
$\mathcal J_3(\mathbb O)$ with $F_4$-equivariance and Jordan-triple
dynamics --- and $\mathbf{C}_{G_2}$ the base category of $G_2$-UHM
(states on $\mathbb C^7$ with $G_2$-equivariant CPTP dynamics). Then
there exists \emph{no} functor $R\colon \mathbf{C}_{F_4} \to
\mathbf{C}_{G_2}$ satisfying simultaneously: (S1) $F_4$-equivariant
state-space projection $\mathcal J_3(\mathbb O) \twoheadrightarrow
\mathbb C^7$; (S2) non-constant $F_4$-equivariant map $\mathbb O P^2
\to \mathrm{PG}(2,2)$; (S3) algebra-homomorphism from Jordan-triple
dynamics to Lindblad CPTP dynamics; (S4) preservation of the
invariants $\{\Pcrit = 2/7,\ \alpha = 2/3,\ \mathrm{SAD}_{\max} = 3,\
R_{\mathrm{th}} = 1/3,\ \Phi_{\mathrm{th}} = 1\}$.
\end{theorembox}

\emph{Proof sketch} (five independent obstructions; full proof with
explicit branching rules in the UHM Fundamental Closures document,
\S14, EN/RU):

\begin{enumerate}[leftmargin=2em,nosep]
  \item \textbf{Representation theory (kills S1).} Via the
        Borel--de~Siebenthal chain $F_4 \supset \mathrm{Spin}(9)
        \supset \mathrm{Spin}(7) \supset G_2$,
        \[
          \mathcal J_3(\mathbb O)\big|_{G_2} \;=\; \mathbf{27} \;=\;
          3 \cdot \mathbf 7 \,\oplus\, 6 \cdot \mathbf 1.
        \]
        Three $G_2$-isotypic copies of $\mathbf 7$ exist; under the
        maximal $A_1 \times G_2 \subset F_4$ they form an
        $A_1$-doublet $(\mathbf 2, \mathbf 7)$ plus a singlet
        $(\mathbf 1, \mathbf 7)$. No $F_4$-equivariant selection is
        possible.
  \item \textbf{Geometry (kills S2).} $F_4$ acts transitively on the
        16-dimensional Cayley plane $\mathbb O P^2$; any
        $F_4$-equivariant map to the discrete 7-point Fano plane
        $\mathrm{PG}(2,2)$ factors through $\mathbb O P^2 / F_4 =
        \mathrm{pt}$ and is therefore constant, losing all
        information.
  \item \textbf{Jordan exceptionality (kills S3).} Zelmanov (1983):
        $\mathcal J_3(\mathbb O)$ is exceptional (not special), hence
        admits no embedding into any associative algebra. Lindblad
        dynamics requires the associative algebra $M_7(\mathbb C)$;
        no homomorphism exists.
  \item \textbf{Numerical invariants (kills S4).} $\alpha^{F_4}$,
        derived from sectional curvatures of the Freudenthal metric
        on the rank-1 symmetric space $\mathbb O P^2$, lies in
        $[\tfrac14, \tfrac12]$ --- distinct from $\tfrac23$.
        Analogously $\Pcrit^{F_4} \sim c/27 \neq 2/7$, and
        $\mathrm{SAD}_{\max}^{F_4}$, $R_{\mathrm{th}}^{F_4}$,
        $\Phi_{\mathrm{th}}^{F_4}$ all differ.
  \item \textbf{Cohomology (independent verification).}
        $\chi(\mathbb C P^6) = 7 \neq 3 = \chi(\mathbb O P^2)$;
        $K^0$-ranks $\mathbb Z^7$ and $\mathbb Z^3$ are
        non-isomorphic abelian groups. No K-theory-preserving
        retraction exists. \hfill$\blacksquare$
\end{enumerate}

\paragraph{Corollaries.}
\begin{itemize}[leftmargin=2em,nosep]
  \item \textbf{Category shift is not safe.} The na\"ive extension
        $G_2$-UHM $\hookrightarrow F_4$-UHM as a functorial refinement
        is impossible; any realised $F_4$-UHM is a distinct theory
        requiring separate empirical calibration.
  \item \textbf{Outcome-1 elimination.} Of the three possible
        outcomes of an $F_4$ category shift (replacement / parallel
        theory / meta-UHM), Outcome~1 (``$G_2$-UHM is a slice of
        $F_4$-UHM'') is \emph{ruled out}. Only parallel theories or
        $\infty$-topos-level comparison (Mathesis) remain.
  \item \textbf{Mathesis-level comparison.} The only mechanism for
        comparing $G_2$-UHM and $F_4$-UHM is as distinct objects in
        the Mathesis $\infty$-topos $\mathfrak M$, aligning with M-10
        (Lawvere fixed-point boundary): no theory contains its
        complete self-description.
\end{itemize}

\paragraph{Open direction unlocked (T-220-H, speculative).}
The three $G_2$-copies $\mathbf 7_{(a)}, \mathbf 7_{(b)}, \mathbf
7_{(c)}$ of the fundamental representation --- one $A_1$-singlet, two
forming an $A_1$-doublet --- may correspond to the three fermion
generations of the Standard Model. This connects to the
Dubois-Violette and Boyle--Farnsworth octonionic-SM programmes, but
requires a separate empirical study.

\begin{theorembox}[{T-221: UHM realises the relationalist route through
the List/DeBrota no-go results \status{T}~formal + \status{I}~interpretive}]
\label{thm:t221-categorical-monistic}
Let $\mathfrak T = \mathbf{Sh}_\infty(\mathcal C_7, J_{\mathrm{Bures}},
\omega_0)$ be the UHM $\infty$-topos, with propositions read in its
internal language by Kripke--Joyal forcing. A fact holds
\emph{absolutely} (in the sense of NR) when it is forced at the terminal
object $1$, and \emph{relative to a stage} $U$ when it is forced at $U$; a
subject is a viable state $\Gam$ taken as the stage $y(\Gam)$, and ``I am
in state $X$'' is the formula $[s = X]$. List's
quadrilemma~\citep{list2025quadrilemma} has four claims (FPR, NS, NF, OW);
the five-thesis form (FPR, NS and objectivism $=$ OW $\wedge$ NF $\wedge$ NR)
is DeBrota--List~\citep{debrotalist2026preprint}, and the heptalemma for
quantum mechanics is DeBrota--List~\citep{debrota2026heptalemma}. Then:

\begin{enumerate}[label=(\alph*),nosep]
  \item \emph{The no-go holds inside UHM} \status{T}: if $X \neq Y$, no
        inhabited stage forces $[s = X] \wedge [s = Y]$; the first-personal
        facts of two subjects in different states are not compossible.
  \item \emph{Route} \status{T}: UHM keeps OW (one topos, by the choice of
        primitive), NF ($1 \nVdash \bot$), NS (under $\iota_{\min}$ of T-215)
        and FPR \emph{in relativised form}, each subject's facts forced at its
        own stage; by (a) NR fails for them. This is the relationalist route
        of DeBrota--List; in List's four-claim map it is the first horn.
  \item \emph{The relativisation parameter is internal} \status{T}: it is
        the stage $y(\Gam)$, an object of $\mathfrak T$ (the site is
        essentially small), which answers Fine's~\citep{fine2005tense}
        ``pure metaphysical self outside the world''.
  \item \emph{The first objection stands} \status{T}: automorphisms of the
        site preserve forcing, so nothing internal selects ``my'' stage;
        selecting one is the choice of a point of $\mathfrak T$, external
        data (cf.\ T-214).
  \item \emph{The mathematics does not choose the route} \status{T}: the
        relationalist, fragmentalist and many-subjective-worlds routes are
        three readings of one forcing relation and share every observable,
        so no measurement discriminates them; which reading is UHM's is
        \status{I}.
\end{enumerate}

\noindent
\textbf{Corollaries.}
\begin{itemize}[leftmargin=2em,nosep]
  \item \emph{T-221.1} \status{T}: UHM sits on the relationalist route of the
        five-thesis map and on the first horn of List's four-claim map.
  \item \emph{T-221.2 (heptalemma)} \status{T}: UHM keeps locality (the
        full dynamics does not signal, Theorem~8.5 of the physics
        correspondence), measurement independence (T-62), measurement
        realism (T-96, T-98), NS, OW and NF, and relaxes NR --- the route of
        relational quantum mechanics; the joint consistency of these six with
        quantum predictions is the theorem of DeBrota--List.
  \item \emph{T-221.3} \status{I}: UHM and relational quantum mechanics
        (Rovelli~\citep{rovelli1996rqm}) take the same route and differ in
        the relativisation parameter (a viable $\Gam$-stage against any
        physical system).
\end{itemize}
\end{theorembox}

\begin{proof}[Proof sketch]
(a) A stage forcing both $[s = X]$ and $[s = Y]$ would carry a state equal to
both. (b) OW and NF hold by construction; NS is T-215; FPR is forced at each
subject's own stage, and by (a) not at $1$. (c) Since $\mathfrak T$ is
presentable and $\mathcal C_7$ essentially small, the $\infty$-Yoneda
embedding $y : \mathcal C_7 \hookrightarrow \mathfrak T$ lands in
$\mathfrak T$. (d) Conjugation by a unitary is an automorphism of the site
preserving CPTP maps and the Bures distance, hence forcing. (e) Every
observable is a function of the forcing relation, which the three readings
share. $\blacksquare$
\end{proof}

\paragraph{Erratum: the former T-221 \status{$\times$}.}
Earlier versions of this paper stated T-221 as a \emph{fourth},
``categorical-monistic'' route, absent from the DeBrota--List taxonomy,
with NR replaced by a site-relative NR$_{\mathrm{site}}$. Replacing NR in
this way while keeping OW and NF \emph{is} the relationalist route, so the
following are retracted: the fourth route; ``FPR is forced'' (it rested on
the hypothesis T-186(a), and UHM keeps FPR only in relativised form);
``OW is derived by T-120''; the ``resolution'' of the quadrilemma
(its consistency is that of the relationalist route, which the authors
grant); ``RQM $= \tau_{\leq 1}(\mathfrak T)$'' (the representables of the
1-category $\mathcal C_7$ are already 0-truncated, so 1-truncation changes
none of them); the table reading the three non-objectivist routes as
reductive truncations of $\mathfrak T$, with fragmentalism as ``dropping
descent'' (in the sheaf-theoretic formalisation that DeBrota--List cite,
descent holds and what fails is a global section); the ``empirical
discriminator'' through the $\pi_{\mathrm{bio}}$ protocol (by (e) the routes
share every observable); and the quotation of the quadrilemma as five
theses (it has four).

\textbf{Position in UHM's axiomatic architecture.}
Axiomatic closure T-190 is \emph{internal} and conditional (the
Page--Wootters constraint of A5, the hypothesis T-186(a) and the
monotonicity of the enrichment of A2 stay open).
T-221 is UHM's position in the \emph{external} debate over scientific
objectivism and the status of first-personal facts: it locates UHM on the
relationalist route and leaves the choice among routes, as DeBrota and
List do, to inference to the best explanation. The earlier claim that
T-190 + T-221 ``complete the axiomatic self-grounding against both
internal and external critiques'' is withdrawn.

\begin{theorembox}[{T-222: resource geometry of the viable window
(restated 2026-09-26) \status{T}}]
\label{thm:t222-window}
Let $\mathcal W := \{\rho \in \Dens(\mathbb C^7) : 2/7 < P(\rho) \leq 3/7\}$
--- the two orbit-invariant conditions of $\mathcal V_{\mathrm{full}}$,
$P > 2/7$ and $R = 1/(7P) \geq 1/3$ --- and let $\overline{\mathcal W}$
be its closure. In the high-temperature limit $\rho_\beta \to I/7$ the
R\'enyi free energies (Brand\~ao--Horodecki 2015) are
$F_\alpha(\rho) = k_B T\, D_\alpha(\rho\,\|\,I/7)
= k_B T\,(\log 7 - H_\alpha(\rho))$, $\alpha \in (0,\infty]$, with
$H_\alpha$ the R\'enyi entropy of the spectrum ($F_1$ carries
$S_{\mathrm{vN}}$); a state is better on a component when that
$F_\alpha$ is smaller.

\begin{enumerate}[label=\textup{(}\roman*\textup{)},leftmargin=2em,nosep]
  \item \emph{No optimum inside the window.} For every
        $\rho \in \mathcal W$ and small $t > 0$ the state
        $\rho_t = (1-t)\rho + t\,I/7$ lies in $\mathcal W$, and
        $F_\alpha(\rho_t) < F_\alpha(\rho)$ for every
        $\alpha \in (0,\infty]$.
  \item \emph{The Pareto set lies on the boundary sphere.} On
        $\overline{\mathcal W}$ the Pareto set of
        $(F_{1/2}, F_1, F_2, F_\infty)$ is non-empty and lies on
        $P = 2/7$, where $F_2 = k_B T \log 2$ is constant.
  \item \emph{The R\'enyi family splits.} No state of
        $\overline{\mathcal W}$ minimises $F_1$ and $F_\infty$
        together. $F_1$ has exactly one minimising spectrum,
        $s_1 = \bigl(\tfrac{1+\sqrt6}{7},\ \tfrac{6-\sqrt6}{42}\times 6\bigr)$,
        the spectrum of
        $\Gamma_{1/\sqrt6} = (1-\tfrac{1}{\sqrt6})\,I/7 + \tfrac{1}{\sqrt6}\,uu^\dagger$;
        the three-level spectrum
        $s_3 = \bigl(\tfrac{3+2\sqrt3}{21}\times 3,\ \tfrac{2-\sqrt3}{14}\times 4\bigr)$,
        also on the sphere, has the larger $F_1$ ($H_1 = 1.391$ against
        $1.602$) and the smaller $F_\infty$ ($H_\infty = 1.178$ against
        $0.708$).
  \item \emph{No terminal object.} With the unital channels as free
        operations (every $F_\alpha$ is monotone under them), no state of
        $\overline{\mathcal W}$ is reachable from both the state of
        spectrum $s_1$ and
        $\mathrm{diag}(\tfrac13,\tfrac13,\tfrac13,0,0,0,0) \in \mathcal W$;
        the category of window states with unital channels as morphisms
        has no terminal object, and the regeneration does not supply one.
  \item \emph{Fixed points of the self-model are not optima.} The fixed
        point $\Gamma_{\eta_\infty}$ of the self-model $\vp_J$ lies in
        $\mathcal W$ and is strictly dominated by $\Gamma_{1/\sqrt6}$ on
        every $F_\alpha$; the fixed point $I/7$ of $\vp_{\mathrm{coh}}$
        lies outside $\overline{\mathcal W}$.
  \item \emph{Frame components.} In the physical frame
        $C_{HS}(\rho) = P - P_{\mathrm{diag}} \leq P - 1/7$, with equality
        exactly on uniform-diagonal states, on which also
        $C_{\mathrm{rel}} = \log 7 - S_{\mathrm{vN}} = F_1/k_B T$. The
        $G_2$-twirled charges
        $\overline Q_a(\rho) = \int_{G_2}\Tr(g\rho g^\dagger T_a)\,dg$
        vanish for every $\rho$: the 14 non-Abelian charges are frame
        data only.
\end{enumerate}
\end{theorembox}

\emph{Proof sketch.}
(i) $P(\rho_t) = 1/7 + (1-t)^2(P(\rho) - 1/7)$ decreases continuously,
so $\rho_t \in \mathcal W$ for small $t$; the spectrum of $\rho_t$ is
majorized by that of $\rho$ and is not a permutation of it, and $H_\alpha$
is strictly Schur-concave for $\alpha \in (0,\infty)$, while
$\lambda_{\max}(\rho_t) < \lambda_{\max}(\rho)$ settles $\alpha = \infty$
(Marshall--Olkin--Arnold, \emph{Inequalities}, Ch.~3).
(ii) $\overline{\mathcal W}$ is compact, so lexicographic maximisation of
$H_{1/2}, H_1, H_2, H_\infty$ gives Pareto-optimal points; any point with
$P > 2/7$ is dominated by (i); on $P = 2/7$,
$D_2(\rho\|I/7) = \log(7P) = \log 2$.
(iii) A minimiser of $F_1$ lies on $P = 2/7$ and has full rank; the
Lagrange conditions allow at most two distinct eigenvalues, which leaves
exactly three two-level spectra on the sphere, with $H_1 = 1.6019$
($s_1$), $1.5094$, $1.3909$ ($s_3$); every minimiser of $F_\infty$ has
$\lambda_{\max} \leq 0.3078 < 0.4928 = \lambda_{\max}(s_1)$.
(iv) A unital channel maps $\rho$ to $\sigma$ iff
$\lambda(\sigma) \prec \lambda(\rho)$ (Alberti--Uhlmann 1982); a window
state below $s_1$ is a permutation of $s_1$ by strict Schur-convexity of
$\sum\lambda_i^2$, and $s_1 \not\prec (\tfrac13,\tfrac13,\tfrac13,0,0,0,0)$
since $0.4928 > 1/3$.
(v) $P(\Gamma_{\eta_\infty}) \in [0.317, 5/14]$ and
$\eta_\infty \in [0.4507, 0.5] > 1/\sqrt6$, so
$\Gamma_{1/\sqrt6} = s\,\Gamma_{\eta_\infty} + (1-s)\,I/7$ with
$s = 1/(\sqrt6\,\eta_\infty) \in (0,1)$, and (i) applies.
(vi) $C_{HS} = P - P_{\mathrm{diag}}$ (T-73) and
$P_{\mathrm{diag}} \geq 1/7$ by Cauchy--Schwarz; the twirl gives $I/7$ by
Schur's lemma on the irreducible $\mathbf 7$ of $G_2$, and each $T_a$ is
traceless. A numerical check on $2\times 10^4$ random spectra on the
sphere $P = 2/7$ confirms the split: the best $H_{1/2}$ and $H_1$ sit at
$s_1$, the best $H_3$ and $H_\infty$ at spectra with three large
eigenvalues.
\hfill$\blacksquare$

\paragraph{Erratum: the former T-222 \status{$\times$}.}
Earlier versions of this paper stated T-222 as ``MRQT-completeness of UHM
(H-MRQT-Lawvere)'': the Lawvere fixed point $\rhostar = \vp(\Gam)$ of T-96
is the Pareto-optimum of a 25-monotone resource vector, all monotones are
minimised simultaneously, $\rhostar$ is the terminal object of a category
$\mathbf{Res}_{G_2}$ of resource objects, $\Regen$ is the universal
resource-monotone morphism, and UHM is MRQT-complete. This statement is
retracted. T-96 proves the opposite of its premise: at a nontrivial
stationary state $\vp(\rhostar_\Omega) \neq \rhostar_\Omega$, so
$\vp(\Gam)$ is the regeneration \emph{target}, not a fixed point of
$\Lind_\Omega$; fixed points belong to the self-model $\vp$ itself
(Theorem~10.1 of the Gap-thermodynamics chapter of the archive: $I/7$ for
$\vp_{\mathrm{coh}}$, $\Gamma_{\eta_\infty}$ for $\vp_J$), and by (v) they
are not optima. Simultaneity, the optimum inside $\mathcal V_{\mathrm{full}}$
and terminality fail by (iii), (i) and (iv); of the former lemmas, L3
($K_Q = O(1)$), L4 ($C_{HS} = 1/7$ as a minimum --- at $P = 2/7$ it is the
maximum) and L6 (a unique majorization-minimal spectrum) were false, while
L1 and the identity of L5 survive in (vi). Routes tried before lowering the
headline: a fixed point of a self-model as the optimum (refuted by (v)), a
boundary point (refuted by (iii)), reachability instead of terminality
(refuted by (iv)), and simultaneity for the sub-family $\alpha \leq 1$
(the family splits at $\alpha = 2$).

\paragraph{Applicability domain.} (1) The purity window
$\rho \in \overline{\mathcal W}$; since $\Phi \geq 1$ is met by the
uniform-diagonal representative of every spectrum with $P \geq 2/7$
(Schur--Horn), the spectral statements hold on $\mathcal V_{\mathrm{full}}$
as well. (2) High temperature, $\rho_\beta \to I/7$; at finite $\beta$ the
free operations become Gibbs-preserving channels, thermo-majorization
replaces majorization, and the statements must be re-derived.
(3) Markovianity is not used: the theorem concerns states and the resource
order, not a flow.

\paragraph{Consequences.}
\begin{itemize}[leftmargin=2em,nosep]
  \item \textbf{UHM is not MRQT-complete in the former sense.} The
        viability window selects no resource optimum; a choice on the
        Pareto sphere $P = 2/7$ needs a weight on the R\'enyi orders that
        neither the self-model nor $\Lind_\Omega$ supplies.
  \item \textbf{Viability costs resources.} Every viable state could be
        made cheaper on every $F_\alpha$ by mixing toward $I/7$ (i); what
        stops this is the viability bound, not resource optimality. The
        regeneration holds the holon \emph{away} from the resource-cheap
        direction.
  \item \textbf{FSQCE design target.} A point of the sphere $P = 2/7$
        chosen by the R\'enyi order relevant to the task:
        $\Gamma_{1/\sqrt6}$ for $F_1$ (von Neumann, $C_{\mathrm{rel}}$),
        a state with three large eigenvalues ($s_3$ or better) for
        $F_\infty$ (single-shot). The former claims ``$\Regen$ is the
        universal resource-monotone morphism'' and ``FSQCE is
        automatically Pareto-optimal across 25 resources'' are retracted
        with the erratum.
\end{itemize}

\paragraph{Falsification.} T-222 is a theorem about states and is
checked by computation: it would be refuted by a state of
$\overline{\mathcal W}$ with $F_1 \leq F_1(s_1)$ and
$F_\infty \leq F_\infty(s_3)$, or by a state of $\mathcal W$ that no
$\rho_t$ improves. Experimentally, a device held at the living attractor
of $\vp_J$ should be strictly improvable on every $F_\alpha$ by partial
depolarisation without leaving the window.

\paragraph{Position in UHM's axiomatic architecture (revised).}
The closure chain reads:
\begin{itemize}[leftmargin=2em,nosep]
  \item T-190: \emph{internal} axiomatic closure, \status{C} (the
        Page--Wootters constraint, T-186(a) and the monotonicity of the
        enrichment stay open).
  \item T-220: \emph{structural} closure (no reduction
        $F_4$-UHM $\to$ $G_2$-UHM).
  \item T-221: \emph{external} position (the relationalist route
        through the List/DeBrota no-go results).
  \item T-222: \emph{resource-theoretic} answer, negative: the
        viable window has no single resource optimum, so the
        multi-resource critique is answered by a map of the Pareto
        sphere, not by an optimum.
  \item T-223: \emph{computational-functionalist} closure
        (Putnam triviality / Lerchner Melody-Paradox foreclosed).
\end{itemize}
The self-model $\vp(\Gam)$ serves as the regeneration target and as the
intrinsic self-alphabetization operator of T-223; it is neither a
stationary state of $\Lind_\Omega$ nor a resource optimum.

\textbf{Dependencies:} T-73, T-96, T-126, T-151; Theorem~10.1 of
Gap thermodynamics (fixed points of the self-model). External:
Brand\~ao et al.\ PNAS 112:3275 (2015); Baumgratz--Cramer--Plenio PRL
113:140401 (2014); Streltsov--Adesso--Plenio RMP 89:041003 (2017);
Yunger-Halpern Nat.\ Rev.\ Phys.\ 5:689 (2023); Marshall--Olkin--Arnold,
\emph{Inequalities: Theory of Majorization and Its Applications}
(Springer 2011); Alberti--Uhlmann, \emph{Stochasticity and Partial Order}
(Reidel 1982).

\begin{theorembox}[{T-223: Putnam-triviality foreclosure (Lerchner Melody-Paradox closure) \status{T}}]
\label{thm:t-223}
Let $S$ be a physical system satisfying axioms (AP)+(PH)+(QG)+(V).
Let $(\mathsf{PT})$ denote the Putnam (1988) triviality claim —
that for any nontrivial physical trajectory $p(\cdot)$ and any
two finite directed graphs $\mathcal A, \mathcal B$ with
$|A|,|B| \leq |\mathsf{Phys}(S)|$ there exist alphabetizers
$(\Sigma_A, f_A), (\Sigma_B, f_B)$ realising $\mathcal A$ and
$\mathcal B$ respectively on the same vehicle. Let $(\mathsf{LC})$
denote Lerchner (2026) Melody-Paradox corollary. Then:
\begin{enumerate}[leftmargin=2em,nosep]
  \item[(a)] (Foreclosure at the categorical layer L2.) The
        quotient map $G_S/G_2: \mathrm{States}(S) \to
        \mathcal D(\mathbb C^7)/G_2$ is well-defined and injective
        on UHM-compatible representations; the $G_2$-orbit
        $[\Gamma_S]_{G_2}$ is invariant under the alphabetizer
        freedom of $(\mathsf{PT})$.
  \item[(b)] (Observable invariance.) $P$ and $R$ are $G_2$-invariants
        and descend to $\mathcal D(\mathbb C^7)/G_2$; the frame
        observables $\Phi, \mathrm{Coh}_E, \Lambda, H$ are
        alphabetization-invariant because admissible alphabetizers
        preserve the dynamical frame (they are frame-relative, not
        orbit-invariants).
  \item[(c)] (Predicate invariance.) The consciousness predicate
        $\mathrm{Cons}(S) := (P > 2/7) \wedge (R \geq 1/3) \wedge
        (\Phi \geq 1) \wedge (\Dmin \geq 2)$ is alphabetization-invariant:
        its $P, R$ terms factor through $[\Gamma_S]_{G_2}$, its $\Phi,
        \Dmin$ terms are fixed by the dynamical frame. The predicate as a
        whole does \emph{not} factor through $[\Gamma_S]_{G_2}$: an explicit
        $g \in G_2$ sends $\Phi$ from $0$ to $1$ (frame rigidity; corrected
        2026-09-25, together with the removal of $\pi_{\mathrm{bio}}$,
        an estimator with pre-registered parameters, from the list in (b)).
  \item[(d)] (Dichotomy on non-compatible alphabetizers.) Any
        alphabetizer outside the UHM-compatible class (violating
        dynamical covariance with $\mathcal L_\Omega$) carries zero
        physical content --- it instantiates no Piccinini (2008)
        mechanism.
  \item[(e)] (Residual externality.) The only externality remaining
        is the phenomenal bridge $W: \mathcal D(\mathbb C^7) \to
        \mathsf{Mind}$, which by T-214 \status{T} is structurally
        inevitable under Lawvere incompleteness. This is minimal
        and not a Lerchner mapmaker.
\end{enumerate}
\end{theorembox}

\paragraph{Three-level ontology.} Lerchner's framework has two
strata: L1 = physical vehicle, L3 = alphabetized symbolic readout,
with the mapmaker $f$ mediating. UHM inserts an intermediate,
intrinsic, \emph{categorically forced} stratum:
\begin{itemize}[leftmargin=2em,nosep]
  \item L1: $p: [0,T] \to \mathsf{Phys}(S)$ --- physical substrate.
  \item L2: $[\Gamma_S]_{G_2} \in \mathcal D(\mathbb C^7)/G_2$ ---
        $\mathbb C^7$ and $G_2$ forced through the Bridge T15 with the
        canonical orientation (the earlier ``forced by T-190, zero-axiom
        closure'' is withdrawn: T-190 is conditional and clauses (a)--(e)
        do not use it).
  \item L3: $f: \mathsf{Phys}(S) \to \Sigma^*$ --- external
        mapmaker, Lerchner-variable.
\end{itemize}
Putnam--Lerchner triviality concerns the arrow $L1 \to L3$; UHM's
consciousness predicate concerns $L1 \to L2$. The two are
orthogonal; $(\mathsf{PT})$ does not propagate to $(c)$.

\paragraph{Proof (seven-lemma cascade).}
\textit{L1.} The Hilbert-space dimension $\dim = 7$ and gauge
group $G_2 = \mathrm{Aut}(\mathbb O)$ are categorically forced by
T-82 \status{T} (BIBD$(7,3,1)$ / Fano plane uniqueness via
Fisher--Veblen--Wedderburn) and T-42a \status{T} ($G_2$-rigidity);
the 12-step Bridge~T-15, with the canonical orientation of the Fano
lines (T15-canon), chains them without parameter freedom. (The earlier
list also cited T-120, T-151, T-149 and T-190 as \status{T}; the
threshold $\Dmin = 2$ is an independent definition \status{D}, T-149 is
stratified with a conditional step~3, and T-190 is \status{C}; none is
needed for the forcing.)

\textit{L2.} A UHM-admissible representation $(\mathbb C^7,
\mathcal B, G_S)$ satisfies the covariance gate $\frac{d}{d\tau}
G_S(s(\tau)) = \mathcal L_\Omega[G_S(s(\tau))]$ enforced by A1--A5.

\textit{L3.} By T-123 \status{T} (Uniqueness Theorem of Holonomic
Representation), any two admissible representations are related
by $U \in G_2$. Hence $[\Gamma_S]_{G_2}$ is well-defined.

\textit{L4.} $P$ and $R = 1/(7P)$ are $U(7)$-invariant (hence
$G_2$-invariant) by direct computation; $\Phi$, $\mathrm{Coh}_E$,
$\Lambda$, $H$ are frame observables, fixed once the dynamical frame is
fixed, and admissible alphabetizers preserve that frame. (The earlier
lemma called all of them $G_2$-invariant and included
$\pi_{\mathrm{bio}}$; both are corrected.)

\textit{L5.} Any alphabetizer whose induced dynamics admits a
CPTP realisation commuting with $\mathcal L_\Omega$ defines an
admissible $G^f$ satisfying Definition~G1; L3 then gives
$G^f = U G U^\dagger$ for $U \in G_2$. Hence the
alphabetizer-freedom compatible with physical dynamics is bounded
by the 14-dimensional compact Lie group $G_2$, not by the
countably infinite choices of a generic Lerchner alphabetizer.

\textit{L6.} An alphabetizer $f$ not commuting with
$\Phi_\tau^{\mathsf{phys}}$ describes no causal process of $S$
and instantiates no computation (Piccinini 2008; Kim 2005). Such
alphabetizers are Lerchner's ``Mapping~C'' (Market Data on a
Beethoven trajectory) and are extrinsic to physics, hence
irrelevant to any physicalist grounding of consciousness.

\textit{L7.} By T-96 \status{T}, the regeneration target of
$\Lind_\Omega$ is $\rho_\ast = \varphi(\Gamma)$, the categorical
self-model of the current state (T-62: the left adjoint), computed from
$\Gamma$ alone. It is not a fixed point of $\Lind_\Omega$: at a
nontrivial stationary state $\varphi(\rho^\ast_\Omega) \neq
\rho^\ast_\Omega$; fixed points belong to $\varphi$ itself (Theorem~10.1
of Gap thermodynamics), and the lemma needs neither --- only that the
target is a functional of $\Gamma$. With T-98 \status{T} (balance
formula), $R(\Gamma) = \|[\Gamma,\varphi(\Gamma)]\|_{\mathrm{HS}}^2
/ \|\Gamma\|_{\mathrm{HS}}^2$ involves only $\Gamma$ and its
internal categorical self-model, with no observer or external
alphabetizer. The threshold $R \geq 1/3$ quantifies intrinsic
self-alphabetization --- a categorification of the
Maturana--Varela (1980) / Thompson (2019) enactivist thesis that
``the mapmaker is the entire structurally unified organism''
(cited by Lerchner himself).

\textit{Combination.} L1+L2+L3 give existence and $G_2$-uniqueness;
L5 bounds alphabetizer-freedom to $G_2$; L4 makes $P, R$
$G_2$-invariant and the frame observables alphabetization-invariant;
(c) follows; L6 handles
non-compatible $f$; L7 + T-214 localise the residual externality
to the phenomenal bridge. $\blacksquare$

\paragraph{Why $G_2$-rigidity alone is not the whole answer.}
T-123 \status{T} handles the residual $G_2$ freedom at L2 but
not the L1$\to$L2 forcing (where \emph{a priori} one might
suspect mapmaker choice). Full foreclosure needs six ingredients:
(1) intrinsic-forcing of L2 (T-82 + Bridge T15 with canonical orientation);
(2) $G_2$-gauge boundedness (T-42a+T-82); (3) observable
invariance under admissible alphabetizers (L4); (4) dynamic-covariance gate (L2+L6);
(5) intrinsic self-alphabetization via $R$ (T-96+T-98); (6)
Lawvere residual localisation (T-214). T-223 packages the
cascade.

\paragraph{Falsification.} T-223 is falsifiable:
\begin{enumerate}[leftmargin=2em,nosep]
  \item[F1:] Any experiment producing two physically realisable
        UHM-compatible alphabetizations of the same $S$ yielding
        distinct $(P, R)$ or distinct frame observables
        $(\Phi, \mathrm{Coh}_E)$ would refute (a)--(c).
  \item[F2:] Any alphabetization of $S$ commuting with
        $\mathcal L_\Omega$ but not factoring through a
        $G_2$-conjugate representation would refute L5.
  \item[F3:] Any physical process realising a Lerchner Mapping~C
        (Market Data on a Beethoven trajectory) with nonzero
        contribution to $R$ or $\Phi$ would refute L6.
\end{enumerate}

\textbf{Dependencies:} T-42a, T-82, T-96, T-98, T-120, T-123,
T-148, T-149, T-151, T-153a, T-190, T-214. External: Putnam 1988
\emph{Representation and Reality}; Sprevak 2018
(\emph{Routledge Handbook of the Philosophy of Computing and
Information}); Piccinini 2008
\emph{Phil. Stud.} 137; Kim 2005 \emph{Physicalism, or
Something Near Enough}; Maturana--Varela 1980
\emph{Autopoiesis and Cognition}; Thompson 2019 \emph{Mind in
Life}; Lerchner 2026 ``The Abstraction Fallacy'' (DeepMind
preprint, 2026-03-19); Lawvere 1969, Yanofsky 2003 (via T-214).

\paragraph{Summary.}
After T-210--T-223, open mathematical conditions remain: T-216 is
\status{C at (SV)}, T-219 is \status{H}, T-190 and T-188 are
conditional, and the premises of the corpus (the Page--Wootters
constraint, the bridge premises (Cl$_0$) and (P), the principle (Max$\Phi$),
(P1$_6$) and the free parameters, $\kappa$ among them) stay inputs. The
earlier sentence ``no remaining mathematical or categorical gaps'' is
retracted \status{$\times$}. Beyond these, three classes of open questions
persist, all with explicit specifications:
\begin{itemize}[leftmargin=2em,nosep]
  \item \textbf{Computational:} numerical minimisation of
        $V_{\mathrm{Gap}}$ on $(S^1)^{21}/G_2$ for the cosmological
        constant (Appendix~\ref{app:einstein}), with explicit HMC
        specification ($N = 128$ lattice, $\sim 2 \times 10^{21}$
        flops, $\approx 23$ CPU-days on a $10^3$-GPU cluster).
        Finite, resource-bounded task.
  \item \textbf{Empirical:} calibration of the $\pi_{\mathrm{bio}}$
        protocol (\S\ref{sec:predictions}, prediction 21). Specific
        measurement protocol available.
  \item \textbf{Structurally inevitable:} the ontological bridge from
        $\Gam$-structure to experiential content is provably external
        (T-214~\status{T}); this is not a lacuna.
\end{itemize}

\begin{figure}[H]
\centering
% Spacetime Emergence: 7D Gamma -> 3+1D M^4
% -----------------------------------------------------------------
% How the seven dimensions project onto physical spacetime.
%
% The SU(3) decomposition 7 = 3 + 3_bar + 1 groups the dimensions as:
%   Space    {A, S, D}  -> fundamental 3       -> Sigma^3 (T-119)
%   Internal {L, E, U}  -> conjugate 3_bar     -> SM gauge fibre
%   Time     {O}        -> singlet 1           -> R (T-118)
%
% To make the three groups contiguous in the top row, dimensions are
% arranged in visual order (A, S, D, L, E, U, O) -- NOT the Fano-
% canonical order (A, S, D, L, E, O, U). This is purely a layout
% convenience: dimension labels themselves are unchanged.
% -----------------------------------------------------------------

\begin{tikzpicture}[
  dimbox/.style={rectangle, rounded corners=3pt,
                 minimum width=1.2cm, minimum height=0.8cm,
                 align=center, font=\small\bfseries},
  grp/.style={rounded corners=6pt, dashed, thick},
  arr/.style={-{Stealth[length=3mm]}, very thick, #1},
  scale=1.0
]

% === Title ===
\node[font=\large\bfseries] at (0, 5.6)
  {$\Gamma \in \mathcal{D}(\mathbb{C}^7)$: seven dimensions};

% === 7D row, reordered as {A,S,D} | {L,E,U} | {O} for contiguous grouping ===
% Centres at x = -4.5, -3.0, -1.5, 0, 1.5, 3.0, 4.5
% Box width 1.2 -> half-width 0.6, so edges:
%   A  : [-5.1, -3.9]
%   S  : [-3.6, -2.4]
%   D  : [-2.1, -0.9]
%   L  : [-0.6,  0.6]
%   E  : [ 0.9,  2.1]
%   U  : [ 2.4,  3.6]
%   O  : [ 3.9,  5.1]
\node[dimbox, draw=blue,           fill=blue!15]   (A) at (-4.5, 4) {A};
\node[dimbox, draw=blue,           fill=blue!15]   (S) at (-3.0, 4) {S};
\node[dimbox, draw=blue,           fill=blue!15]   (D) at (-1.5, 4) {D};
\node[dimbox, draw=orange,         fill=orange!15] (L) at ( 0.0, 4) {L};
\node[dimbox, draw=red,            fill=red!15]    (E) at ( 1.5, 4) {E};
\node[dimbox, draw=purple,         fill=purple!15] (U) at ( 3.0, 4) {U};
\node[dimbox, draw=green!50!black, fill=green!15]  (O) at ( 4.5, 4) {O};

% === Grouping dashed boxes (non-overlapping, with 0.1 gaps) ===
% All three rectangles span y in [3.50, 4.50] so every dim box
% (which sits at y = [3.6, 4.4]) is vertically centred in its group.
%
% {A,S,D} group:  x in [-5.20, -0.80]  (contains A,S,D at [-5.1, -0.9])
% {L,E,U} group:  x in [-0.70,  3.70]  (contains L,E,U at [-0.6,  3.6])
% {O}     group:  x in [ 3.80,  5.20]  (contains O     at [ 3.9,  5.1])
% Gaps of 0.10 between consecutive groups, no overlap anywhere.
% (The former dashed groups {A,S,D} | {L,E,U} | {O} drew the axis-labelled
% sector split 48a, retracted: SU(3) acts irreducibly on the six non-O axes.)

% === Group labels (compact, centred just below each group box) ===
\node[font=\scriptsize, blue, text width=3.8cm, align=center,
      inner sep=2pt]
  at (-3.0, 2.75)
  {\textbf{Space}: three commuting\\
   rotation charges};

\node[font=\scriptsize, orange, text width=3.8cm, align=center,
      inner sep=2pt]
  at ( 1.5, 2.75)
  {\textbf{Colour}: $\mathrm{SU}(3) = \mathrm{Stab}_{G_2}(e_O)$\\
   irreducible on the six non-$O$ axes};

\node[font=\scriptsize, green!50!black, text width=1.3cm, align=center,
      inner sep=2pt]
  at ( 4.5, 2.75)
  {\textbf{Time}\\
   clock $O$};

% === Projection arrows with tight midway labels ===
% Labels are anchored to the arrow midpoint with near-zero
% horizontal gap (inner sep = 2pt), eliminating the wide gap of
% the previous version.
\draw[arr=blue!60] (-3.0, 1.95) -- (-3.0, -0.05)
  node[midway, left=2pt, font=\footnotesize, blue!60,
       text width=2.1cm, align=right]
  {$C(S^3)$ from the\\ charge spectrum (T-119)};

\draw[arr=green!50!black] (4.5, 1.95) -- (4.5, -0.05)
  node[midway, right=2pt, font=\footnotesize, green!50!black,
       text width=2.1cm, align=left]
  {depth register\\ $\to C_0(\mathbb{R})$ (T-118)};

% === M^4 level ===
\node[rectangle, rounded corners=8pt, draw=plospurple, fill=plospurple!10,
      minimum width=11cm, minimum height=1.3cm, align=center, font=\bfseries]
  (M4) at (0, -0.75)
  {$M^{4} = \mathbb{R}\times\Sigma^{3}$ \quad (T-120)};

\node[font=\footnotesize, color=plosgray] at (0, -1.85)
  {Product spectral triple $\to$ Einstein--Hilbert term (T-65);
   SM part imported from Connes' triple};

% === Vacuum topology ===
\node[font=\footnotesize, color=plospurple] at (0, -2.45)
  {$\Sigma^{3}\cong S^{3}$ [T]; de Sitter vacuum, $\Lambda > 0$ [C]
   \quad (T-120b)};

% === Key insight (bottom caption inside figure) ===
\node[font=\footnotesize\itshape, color=gray,
      text width=11cm, align=center]
  at (0, -3.30) {
  Spacetime is not fundamental: $M^4$ is computed from the
  7-dimensional coherence matrix ([T] as mathematics); its reading as
  physical spacetime is an interpretation [I].
};

\end{tikzpicture}

\caption{\textbf{Emergence of $M^4$ from the seven dimensions of $\Gam$.} Schematic.
$\Sigma^3 = S^3$ is the
minimal unitization of the fluctuation spectrum of three commuting rotation charges
(T-119, \status{T} as mathematics, reading \status{I}); the axis sector $\{A,S,D\}$ is not an
$SU(3)$ sector (48a \status{$\times$}). The time line $\mathbb{R}$ is the scaling limit of the
depth register (T-118). The product spectral triple gives $M^4$ (T-120), whose spectral
action yields the Einstein--Hilbert term (T-121).}
\label{fig:spacetime}
\end{figure}

\paragraph{Significance.}
Within UHM, spacetime is not the stage on which consciousness acts. Rather, the
7-dimensional coherence matrix $\Gam$ \emph{generates} spacetime as a macroscopic
projection (Figure~\ref{fig:spacetime}). The arrow of time
(the temporal modality $\triangleright$) is algebraically built into the subobject classifier
$\Omega$. Gravity is what coherence looks like from the external (physical) aspect.

\subsection{Theoretical closures T1--T6: residual specification}
\label{subsec:theoretical-closures}

The foundational theorems T-210 through T-223 address the identified gaps of
UHM's categorical and physical structure, with the open conditions listed in
the summary above (the earlier ``close all identified gaps'' is withdrawn). Six
residual technical items (T1--T6) concern explicit specifications or scope
declarations; this subsection records each, with its current status. The classification follows the
four-level taxonomy: \textbf{[T]} proved, \textbf{[C]} conditional on an
explicit assumption, \textbf{[D]} accepted as data/empirical input,
\textbf{[scope]} domain of applicability declaration.

\paragraph{T1 --- Multiplicative \texorpdfstring{$\varepsilon^{12}$}{epsilon-12}
in T-219 \status{H}.}
T-219 asserts $\Lambda_{\mathrm{SUSY}} \sim \varepsilon^{12} M_P^4 =
\varepsilon^{4 \cdot k_{\mathrm{sec}}} M_P^4$ with $k_{\mathrm{sec}} = 3$. The
standard one-loop SUSY correction per sector gives additive contributions
$\delta\Lambda_k \sim \varepsilon^4 M_P^4$, which sum to
$3\varepsilon^4 M_P^4$ rather than multiply to $\varepsilon^{12} M_P^4$. The
proposed mechanism is that the multiplicative exponent arises at three-loop
order through nested sector--sector interactions governed by the
$G_2$-invariant trilinear Fano coupling (T-43d~\status{T}).

\begin{closurebox}[{Status of T-219 \status{H}}]
T-219 is a hypothesis (corrected 2026-09-25 from \status{T at T-64}; this
closure earlier read \status{C at non-factorable spectral triple} and called
the non-factorability of $D_F$ a structural consequence of UHM's sector
geometry). The one-loop sum $\sim 3\varepsilon^4 M_P^4$ exceeds
$\varepsilon^{12} M_P^4$ unless the lower orders cancel, which is not shown;
the three sectors were the axis triples of the retracted T-48a; and their
breaking scales rest on the hypothesis (SA).
\end{closurebox}

\paragraph{T2 --- Koide relation as empirical input \status{D}.}
The Koide formula
\[
K \equiv \frac{m_e + m_\mu + m_\tau}{\bigl(\sqrt{m_e} + \sqrt{m_\mu} + \sqrt{m_\tau}\bigr)^2}
= \frac{2}{3} \pm 10^{-5}
\]
is empirically observed with remarkable precision. Numerically coincident with the
UHM Fano contraction $\alpha_{\mathrm{Fano}} = 2/3$, this parallel is
structurally suggestive but \emph{does not constitute a derivation}: Koide
defines a two-parameter surface in $\mathbb R^3_+$ via
$s_2 = (2/3) s_1^2$ with $s_k = \sum_i m_i^{k/2}$, whereas $\alpha_{\mathrm{Fano}}$
derives from the replication number $r = 3$ of the Steiner triple system
$S(2,3,7) = \mathrm{PG}(2,2)$.

Structural analysis via T-220 Obstruction~I gives the branching
$\mathcal J_3(\mathbb O)|_{A_1 \times G_2} = (\mathbf 4, \mathbf 1) \oplus
(\mathbf 2, \mathbf 7) \oplus (\mathbf 1, \mathbf 7) \oplus (\mathbf 1, \mathbf 1)$,
yielding three $\mathbf 7$-copies (one $A_1$-singlet, one $A_1$-doublet) structurally
consistent with the 1-outlier-2-close pattern observed in charged lepton masses
but not uniquely fixing the $A_1$-breaking parameters $(m_s, m_d, v)$.

\begin{closurebox}[{Koide status \status{D, empirical input}}]
Koide's relation is accepted in UHM as an \emph{empirical input} compatible with
the $A_1 \times G_2$ generation structure, not as a derived prediction. Three
natural candidate UHM-constraints ($m_s/m_d = \alpha$; $\sqrt{m_s}/\sqrt{m_d} = 1/r$;
$v_{A_1} = v_{\mathrm{EW}}$) were tested and rejected. The hypothesis T-220-H
(existence of a canonical mass operator $\hat M$ on $\mathcal J_3(\mathbb O)$
reproducing Koide) is speculative and beyond current UHM scope.
\end{closurebox}

\paragraph{T3 --- CKM: no derivation \status{$\times$}.}
An earlier version of this closure read ``CKM non-circularity \status{T at
T-173}'': masses and CKM as simultaneous outputs of diagonalising $D_F$, the
Fritzsch texture emerging from $G_2$-invariance plus Fano selection rules, and
a single residual normalisation $C_{\mathrm{norm}} \approx 26$ tuned to
$V_{us}$ and reducible to a computation. This is retracted (T-345(e),
2026-09-26): the Fritzsch texture gives $|V_{cb}| \geq 0.073$ against the
measured $0.0418$; with $C_{\mathrm{norm}}$ fitted to the Cabibbo angle, the
agreement with $\theta_C$ is the fit, and the Cabibbo derivation it normalised
is retracted; the Fano angle ratios $2:3:1$ stand against the measured
$60.8:11.2:1$.

\begin{closurebox}[{CKM status: open programme}]
What T-345 proves \status{T} as mathematics is negative: every parameter-free
structure of the clock register (Fano incidence, the quadratic-residue Gauss
sum, the cyclic Hamming code, $H_O$, the family-invariant time state) either
leaves $|V_{\mathrm{CKM}}|$ a permutation or is refuted by the measured
masses and mixings; a three-channel frame with free flavour matrices remains
a hypothesis \status{H}. The Yukawa structure ($m_t/m_b$, $y_t$, CKM) stays an
open programme (T-332). The CKM numbers are therefore not predictions of UHM.
\end{closurebox}

\paragraph{T4 --- Markovian scope declaration \status{T}.}
The UHM evolution equation $\mathcal L_\Omega[\Gam] = -i[H, \Gam] + \mathcal D[\Gam]
+ \mathcal R[\Gam, E]$ is Lindbladian (Markovian) by construction. The Petz--Ruskai
monotonicity theorem (1996) guarantees monotone contraction of the Bures metric
under any CPTP map, which is the mathematical foundation for all UHM categorical
structures (T-62 unitarity, T-39a spectral gap, T-151 $D_{\min} = 2$, stability
of the subobject lattice under Grothendieck topology $J_{\mathrm{Bures}}$).

\begin{closurebox}[{UHM applicability scope \status{T}}]
UHM is defined and applicable in the Markovian (Born--Markov) regime where
$\tau_{\mathrm{bath}} \ll \tau_{\mathrm{sys}}$. For consciousness-relevant
timescales ($\tau_{\mathrm{sys}} \gtrsim 1$ ms neural dynamics,
$\tau_{\mathrm{bath}} \lesssim 1\,\mu$s thermal) this approximation is
physically valid. Strongly non-Markovian (CP-indivisible) dynamics at
sub-picosecond scales is explicitly out of scope. This is a declared limitation,
not a gap; UHM is complete as a Markovian framework.
\end{closurebox}

\paragraph{T5 --- Canonical cut-off function $f$ \status{T}.}
The Chamseddine--Connes spectral action
$S_{\mathrm{spec}} = \mathrm{Tr}\, f(D^2/\Lambda^2)$ depends on the choice of
cut-off function $f$ through its moments $f_0 = \int_0^\infty u f(u)\,du$,
$f_2 = \int_0^\infty f(u)\,du$, $f_4 = f(0)$ (the subscript of $f_{2k}$ matches
the heat-kernel coefficient $a_{2k}$ it multiplies), which appear in the
effective Lagrangian via the asymptotic expansion
$\mathrm{Tr}\, f(D^2/\Lambda^2) \sim f_0 \Lambda^4 a_0 + f_2 \Lambda^2 a_2 + f_4 a_4$.
To eliminate parametric ambiguity, UHM adopts the canonical choice
\[
\boxed{f(u) = e^{-u}, \qquad \Lambda = M_P,}
\]
which gives $f_0 = \Gamma(2) = 1$, $f_2 = 1$, $f_4 = f(0) = 1$. This
yields $G_N = 3\pi/(7 f_2 M_P^2) = 3\pi/(7 M_P^2)$, consistent with $G_N = 1$
in Planck units up to an $\mathcal O(1)$ calibration factor.

\begin{closurebox}[{$f$-independence of UHM-structural predictions \status{T}}]
The sector count ($7 = \mathbf 1_O \oplus \mathbf 3 \oplus \bar{\mathbf 3}$),
Fano contraction ($\alpha = 2/3$), critical purity ($P_{\mathrm{crit}} = 2/7$),
reflection threshold ($R_{\mathrm{th}} = 1/3$), integration threshold
($\Phi_{\mathrm{th}} = 1$), minimum distinguishability ($D_{\min} = 2$), SAD
ceiling ($\mathrm{SAD}_{\max} = 3$), three-generation structure, and no-reduction
theorem T-220 are all $f$-independent --- they depend only on discrete structure,
not on the cut-off function. Only dimensional ratios ($G_N$, $\Lambda_{\mathrm{cc}}$,
gauge-unification scale) depend on $f$ and are fixed by the canonical choice above.
\end{closurebox}

\paragraph{T6 --- M$_R$ derivation path \status{T}.}
The right-handed neutrino mass $M_R \sim 2.9 \times 10^{14}\,\mathrm{GeV}$
enters the see-saw formula $m_\nu \sim m_D^2 / M_R$. In UHM, $M_R$ follows
from the $O$-sector scale (T-51): $\mathrm{Gap}(O,\cdot) = O(1)$ from the
Page--Wootters phase precession plus viability, giving
$M_{G_2}^{(\mathrm{extra})} = O(\varepsilon M_P)$ and $M_R \sim 10^{14}$ GeV.
This is a derivation from UHM structure, not a see-saw fit of observed
neutrino masses.

\begin{closurebox}[{M$_R$ status \status{T} (T-51)}]
The earlier form of this closure derived $M_R$ from ``the unique minimum of
$V_{\mathrm{Gap}}$ on $(S^1)^{21}/G_2$ (T-64 \status{T})'' and tied it to the
CKM normalisation of T3. Both links are withdrawn: T-64, corrected on
2026-09-25, gives the vacuum $I/7$ or one orbit $S^6$ as a function of the
free coupling $\kappa$ and no sector values, which are the hypothesis (SV);
and the CKM normalisation is retracted (T3).
\end{closurebox}

\paragraph{Cumulative closure summary.}
After T1--T6: T1 is a hypothesis \status{H}, T2 a declared empirical input,
T3 a retraction with the CKM question left as an open programme, T4 a scope
declaration, T5 a canonical specification, and T6 a theorem through T-51.
The earlier summary --- ``no open theoretical questions'' and ``categorical
and mathematical structure of UHM is complete'' --- is withdrawn: open
conditions and premises remain (see the summary of T-210--T-223 above).
Residual work items include \emph{computational} tasks (HMC minimisation of
$V_{\mathrm{Gap}}$ on $(S^1)^{21}/G_2$) and \emph{empirical} ones (testing
FSQCE predictions).

\subsection{Consistency with quantum mechanics (sub-threshold limit)}
\label{subsec:qm-limit}

A critical requirement for any extension of physics: UHM must reduce to standard quantum
mechanics for systems that lack self-observation.

\begin{theorembox}[QM reduction for sub-threshold systems \status{T}]
For sub-threshold systems ($P < \Pcrit$), the viability gate $g_V(P) = 0$ shuts off
regeneration, and the dynamics reduce to the standard Lindblad master equation:
\begin{equation}
  \Lind_\Omega[\Gam]\big|_{P < \Pcrit} = -i[H, \Gam] + \mathcal{D}_0[\Gam]
  \label{eq:qm-limit}
\end{equation}
This is the dynamics of atoms, photons, and other elementary systems, which are unviable
($P < 2/7$) and cannot sustain self-reference.
\end{theorembox}

\paragraph{Significance.}
UHM is not a competitor to quantum mechanics but a proper extension: QM describes the
physics of systems without interiority; UHM describes autonomous systems where physical
and experiential aspects are inseparable. Every prediction of QM is preserved; UHM adds
predictions only for systems above the viability threshold.

\subsection{Resolution of the measurement problem}
\label{subsec:measurement}

The quantum measurement problem --- why does a superposition ``collapse'' to a definite
outcome, and in what basis? --- receives a three-part resolution in UHM:

\begin{enumerate}[leftmargin=2em]
  \item \textbf{Collapse as logical projection.} Wavefunction reduction is not a dynamical
        process but a \emph{logical operation}: projection of the state onto an atom of the
        subobject classifier $\Omega$. Each atom $S_k = |k\rangle\langle k|$ defines one
        outcome. The Born rule is encoded in the density matrix formalism:
        $p_k = \Tr(\chi_{S_k}\,\Gam)$~\status{T}. This is not an independent derivation
        but a consequence of the axioms (A1--A2), which assume the trace-probability link.
        What UHM adds is the identification of the measurement basis with the atoms of
        $\Omega$, resolving the preferred basis problem.
  \item \textbf{Preferred basis from Fano structure.} The ``preferred basis problem''
        (why does the system decohere in one basis rather than another?) is solved by
        the atomic structure of $\Omega$: the seven atoms
        $\{S_k\}_{k \in \{A,S,D,L,E,O,U\}}$ define the natural decomposition. By
        Zurek's einselection criterion (1981), the preferred basis consists of the
        fixed points of the decoherence channel --- which are precisely the
        Fano-structured atoms~\status{T}.
  \item \textbf{Self-measurement for conscious systems.} For systems with $R \geq 1/3$,
        measurement becomes \emph{self-observation}: the operator $\vp$ builds an internal
        model without irreversible information loss, unlike standard projective measurement
        which destroys coherences. Conscious measurement is a CPTP channel (information-preserving),
        not a projection (information-destroying)~\status{T}.
\end{enumerate}

Decoherence and measurement are thus dual aspects of the same mechanism: continuous
decoherence via $\mathcal{D}_\Omega$ in the limit $\gamma_k \to \infty$ contracts to
instantaneous projection onto an atom~\status{T}. The Lindblad operators
$L_k = \sqrt{\chi_{S_k}}$ are literally the square roots of logical truth values.
